\chapter{Lifting of Quantum Torsion}
\label{chap:torsion}

\section{Jet Prolongation}
\label{sec:jet-prolongation}

\begin{proposition}[Jet Prolongation]
  \label{prop:jet-prolongation}
  \uses{def:quantum-observable-universe}
Let $(\mathcal{C}, A, \omega)$ be a quantum observable universe over a
ring $R$ (\hyperref[def:quantum-observable-universe]{Definition}).  Define the
\emph{jet object}
\[
J^1(A) \;:=\; A \otimes \Omega^1(A),
\]
where $\Omega^1(A) = \ker(m \colon A \otimes A \to A)$ is the bimodule of
Kähler differentials.  Then $J^1(A)$ carries a canonical Frobenius--Hopf
structure inherited from $A$, and the quadruple
$(\mathcal{C},\, J^1(A),\, \omega',\, q)$ --- where
$\omega' = \omega \circ (\text{--} \otimes \Omega^1)$ is the induced fibre
functor --- is again a quantum observable universe over $R$.

---

The proof proceeds in the rigid monoidal category underlying the quantum
observable universe.  Let $A$ denote the Frobenius-Hopf object
representing the gauge-bundle data of $(M,G,P,E)$.

\emph{Step 1: Categorical jet object.}
The first jet bundle $J^1 E$ is represented categorically by the object
$A \otimes \Omega^1(A)$, where
$\Omega^1(A) = \ker(m\colon A\otimes A \to A)$
is the kernel of the multiplication (the bimodule of Kähler
differentials, i.e. the universal first-order differential calculus
over $A$).
Because $A$ is Frobenius, the non-degenerate pairing
$\varepsilon \circ m \colon A \otimes A \to \mathbf{1}$ exhibits
$\Omega^1(A)$ as a direct summand of $A \otimes A$, hence
$J^1(A) := A \otimes \Omega^1(A)$ is again an object of the ambient
category.

\emph{Step 2: Inherited Frobenius-Hopf structure.}
The Frobenius relation
$(\mathrm{Id} \otimes m) \circ (\Delta \otimes \mathrm{Id})
 = \Delta \circ m
 = (m \otimes \mathrm{Id}) \circ (\mathrm{Id} \otimes \Delta)$
ensures that the comultiplication $\Delta\colon A \to A \otimes A$ restricts
to a well-defined comultiplication
$\Delta'\colon J^1(A) \to J^1(A) \otimes J^1(A)$, and the counit
$\varepsilon\colon A \to \mathbf{1}$ induces a counit $\varepsilon'$ on
$J^1(A)$.  The Hopf antipode $S$ extends to $J^1(A)$ by functoriality of the
Kähler differential construction.  Thus $J^1(A)$ is a Frobenius-Hopf object.

This step relies on the \emph{strong multiplicativity} condition: the
co-multiplication must be compatible with the deformed product at every
filtration degree.  Over fields of positive characteristic, the
rigorous proof that this compatibility survives the deformation is the
main theorem of [fred2005], which constructs $\mathrm{Dist}(G)$
as a Gröbner--Shirshov deformation of the commutative
$\mathrm{Dist}(H)$ and verifies the Frobenius relation at each stage
(see also \hyperref[prop:skew-poisson-integration]{Proposition} in
\hyperref[chap:p-adic]{Chapter}).

\emph{Step 3: Non-degeneracy and vanishing.}
The non-degeneracy of the Frobenius form on $A$ (the pairing
$\varepsilon \circ m$ is a perfect pairing) implies non-degeneracy of the
induced form on $J^1(A)$, since $\Omega^1(A)$ is a direct summand.
For the vanishing condition, $\mathrm{Hom}(\mathbf{1}, \mathrm{Ad}(A)) = 0$
implies $\mathrm{Hom}(\mathbf{1}, \mathrm{Ad}(J^1(A))) = 0$: the adjoint
representation of $J^1(A)$ contains $\mathrm{Ad}(A)$ as a subobject, and any
invariant in $\mathrm{Ad}(J^1(A))$ would project to a non-zero invariant in
$\mathrm{Ad}(A)$, contradicting the hypothesis.

\emph{Step 4: Geometric realisation.}
Applying the fibre functor $\omega\colon \mathcal{C} \to \mathrm{Vect}_k$
(where $\mathcal{C}$ is the rigid monoidal category of the quantum
observable universe)
recovers the geometric data: $\omega(J^1(A))$ is the space of sections of the
first jet bundle $J^1 E$, the frame bundle $P'$ and orthonormal frame bundle
$E'$ are the geometric counterparts of the Frobenius-Hopf structure on
$J^1(A)$, and the connection $A'$ is the image under $\omega$ of the
comultiplication $\Delta'$.  Its uniqueness follows from the irreducibility
of $A'$, which is guaranteed by the inherited vanishing condition of
Step 3.
\end{proposition}
\begin{proof}
The proof proceeds in the rigid monoidal category underlying the quantum
observable universe.  Let $A$ denote the Frobenius-Hopf object
representing the gauge-bundle data of $(M,G,P,E)$.

\emph{Step 1: Categorical jet object.}
The first jet bundle $J^1 E$ is represented categorically by the object
$A \otimes \Omega^1(A)$, where
$\Omega^1(A) = \ker(m\colon A\otimes A \to A)$
is the kernel of the multiplication (the bimodule of Kähler
differentials, i.e. the universal first-order differential calculus
over $A$).
Because $A$ is Frobenius, the non-degenerate pairing
$\varepsilon \circ m \colon A \otimes A \to \mathbf{1}$ exhibits
$\Omega^1(A)$ as a direct summand of $A \otimes A$, hence
$J^1(A) := A \otimes \Omega^1(A)$ is again an object of the ambient
category.

\emph{Step 2: Inherited Frobenius-Hopf structure.}
The Frobenius relation
$(\mathrm{Id} \otimes m) \circ (\Delta \otimes \mathrm{Id})
 = \Delta \circ m
 = (m \otimes \mathrm{Id}) \circ (\mathrm{Id} \otimes \Delta)$
ensures that the comultiplication $\Delta\colon A \to A \otimes A$ restricts
to a well-defined comultiplication
$\Delta'\colon J^1(A) \to J^1(A) \otimes J^1(A)$, and the counit
$\varepsilon\colon A \to \mathbf{1}$ induces a counit $\varepsilon'$ on
$J^1(A)$.  The Hopf antipode $S$ extends to $J^1(A)$ by functoriality of the
Kähler differential construction.  Thus $J^1(A)$ is a Frobenius-Hopf object.

This step relies on the \emph{strong multiplicativity} condition: the
co-multiplication must be compatible with the deformed product at every
filtration degree.  Over fields of positive characteristic, the
rigorous proof that this compatibility survives the deformation is the
main theorem of [fred2005], which constructs $\mathrm{Dist}(G)$
as a Gröbner--Shirshov deformation of the commutative
$\mathrm{Dist}(H)$ and verifies the Frobenius relation at each stage
(see also \hyperref[prop:skew-poisson-integration]{Proposition} in
\hyperref[chap:p-adic]{Chapter}).

\emph{Step 3: Non-degeneracy and vanishing.}
The non-degeneracy of the Frobenius form on $A$ (the pairing
$\varepsilon \circ m$ is a perfect pairing) implies non-degeneracy of the
induced form on $J^1(A)$, since $\Omega^1(A)$ is a direct summand.
For the vanishing condition, $\mathrm{Hom}(\mathbf{1}, \mathrm{Ad}(A)) = 0$
implies $\mathrm{Hom}(\mathbf{1}, \mathrm{Ad}(J^1(A))) = 0$: the adjoint
representation of $J^1(A)$ contains $\mathrm{Ad}(A)$ as a subobject, and any
invariant in $\mathrm{Ad}(J^1(A))$ would project to a non-zero invariant in
$\mathrm{Ad}(A)$, contradicting the hypothesis.

\emph{Step 4: Geometric realisation.}
Applying the fibre functor $\omega\colon \mathcal{C} \to \mathrm{Vect}_k$
(where $\mathcal{C}$ is the rigid monoidal category of the quantum
observable universe)
recovers the geometric data: $\omega(J^1(A))$ is the space of sections of the
first jet bundle $J^1 E$, the frame bundle $P'$ and orthonormal frame bundle
$E'$ are the geometric counterparts of the Frobenius-Hopf structure on
$J^1(A)$, and the connection $A'$ is the image under $\omega$ of the
comultiplication $\Delta'$.  Its uniqueness follows from the irreducibility
of $A'$, which is guaranteed by the inherited vanishing condition of
Step 3.
\end{proof}

\begin{corollary}[Archimedean realisation of jet prolongation]
  \label{cor:jet-prolongation-C}
  \uses{prop:jet-prolongation, def:c-quantum-observable-universe}
Let $(M, G, P, E)$ be an archimedean quantum observable universe
(\hyperref[def:c-quantum-observable-universe]{Example}) with instanton connection
$A$ on $P$ and metric $g$ on $E$.  Write $\widetilde{M} = E$ for the total
space of $E$, define $P'$ as the frame bundle and $E'$ as the bundle of
orthonormal frames of the first jet bundle $J^1 E$.  Then $A$ and $g$
uniquely determine a connection $A'$ on $P'$ such that
$(\widetilde{M}, G, P', E')$ is an archimedean quantum observable universe.

---

Apply \hyperref[prop:jet-prolongation]{Proposition} to the categorical data
$(\mathcal{C}, A, \omega)$ underlying the archimedean quantum
observable universe (\hyperref[def:c-quantum-observable-universe]{Example});
let $q$ denote the substrate parameter,
where $\mathcal{C} = \mathrm{Rep}(\mathrm{SU}_{q}(2))$,
$A = \mathcal{O}_q(\mathrm{SU}(2))$, and $\omega$ is the forgetful
functor.  The fibre functor $\omega$ converts the categorical jet object
$J^1(A)$ into the geometric first jet bundle $J^1 E$, the inherited
Frobenius--Hopf structure into the connection $A'$, and the
non-degeneracy and vanishing conditions into the analytic conditions of
\hyperref[def:c-quantum-observable-universe]{Example}.
\end{corollary}
\begin{proof}
Apply \hyperref[prop:jet-prolongation]{Proposition} to the categorical data
$(\mathcal{C}, A, \omega)$ underlying the archimedean quantum
observable universe (\hyperref[def:c-quantum-observable-universe]{Example});
let $q$ denote the substrate parameter,
where $\mathcal{C} = \mathrm{Rep}(\mathrm{SU}_{q}(2))$,
$A = \mathcal{O}_q(\mathrm{SU}(2))$, and $\omega$ is the forgetful
functor.  The fibre functor $\omega$ converts the categorical jet object
$J^1(A)$ into the geometric first jet bundle $J^1 E$, the inherited
Frobenius--Hopf structure into the connection $A'$, and the
non-degeneracy and vanishing conditions into the analytic conditions of
\hyperref[def:c-quantum-observable-universe]{Example}.
\end{proof}

\begin{remark}[Prolongation from the Frobenius-Hopf structure]
  \label{rem:prolongation-categorical}
  \uses{prop:jet-prolongation, def:recursive-brane-tower}
The proof above shows that jet prolongation is an \emph{intrinsic}
consequence of the Frobenius-Hopf structure in a rigid monoidal category:
no additional geometric input beyond the categorical data
$(\mathcal{C}, A, \omega)$ is required.  Concretely:

\begin{enumerate}
  \item The Frobenius structure provides the Kähler differentials
\end{enumerate}
   $\Omega^1(A)$ as a direct summand, from which the jet object
   $J^1(A)$ is built.
\begin{enumerate}
  \item The Frobenius relation ensures that the comultiplication and
\end{enumerate}
   multiplication are compatible, so the prolonged object inherits
   the full Hopf algebra structure.
\begin{enumerate}
  \item The non-degeneracy and vanishing conditions propagate from $A$ to
\end{enumerate}
   $J^1(A)$ by functoriality.

In this sense, the recursive brane tower
(\hyperref[def:recursive-brane-tower]{Definition}) is generated entirely by
iterating the categorical jet construction.  The tower, together with its
associated descent (\hyperref[ssec:descent]{Subsection}), is a formal consequence
of the axioms of a quantum observable universe and does not require any
case-by-case verification.
\end{remark}

\section{Recursive Brane Tower}
\label{sec:recursive-brane-tower}

% ERROR: block "recursive-brane-tower" not found
\begin{example}[Archimedean realisation of the brane tower]
  \label{ex:brane-tower-levels}
  \uses{def:c-quantum-observable-universe, def:recursive-brane-tower, cor:jet-prolongation-C}
Starting from an archimedean quantum observable universe $(M, G, P, E_1)$
(\hyperref[def:c-quantum-observable-universe]{Example}), the fibre functor
$\omega$ converts the categorical brane tower into a geometric tower of
bundles: $E_{n+1}$ is the bundle of orthonormal frames of $J^1 E_n$, and
the connection $A_n$ on $E_n$ satisfies the Hermitian Yang--Mills equations
at each level.  The first jet prolongation yields $E_2$ with base manifold
$M_1 = E_1$, and the connection $A_2$ is uniquely determined by
\hyperref[cor:jet-prolongation-C]{Corollary}.
\end{example}

\section{Yekutieli Lifting and Quantum Characteristic Classes}
\label{sec:yekutieli-lifting}

% ERROR: block "yekutieli-lifting-condition" not found
% ERROR: block "quantum-characteristic-classes" not found
\begin{remark}[Archimedean band of the gerbe extension]
  \label{rem:yekutieli-archimedean-band}
  \uses{def:c-quantum-observable-universe}
In the archimedean realisation
(\hyperref[def:c-quantum-observable-universe]{Example}), the banding sheaf is
$\mathcal{B} = \mathcal{O}_M^{\times}$ (the sheaf of invertible
holomorphic functions on $M$).
\end{remark}

\begin{proposition}[Yekutieli Lifting Condition holds in the brane tower]
  \label{prop:yekutieli-lifting}
  \uses{def:quantum-observable-universe, def:recursive-brane-tower, def:yekutieli-lifting-condition, def:quantum-characteristic-classes, prop:jet-prolongation}
Let $(\mathcal{C}, A, \omega)$ be a quantum observable universe
(\hyperref[def:quantum-observable-universe]{Definition}) and let
$\{A_n\}_{n \geq 1}$ be the recursive brane tower
(\hyperref[def:recursive-brane-tower]{Definition}).  The central extensions of
gerbes
\[
1 \;\longrightarrow\; \mathcal{B}
  \;\longrightarrow\; \mathfrak{G}_{n+1}
  \;\longrightarrow\; \mathfrak{G}_n
  \;\longrightarrow\; 1
\]
arising from the jet prolongation
(\hyperref[prop:jet-prolongation]{Proposition}) satisfy the Yekutieli Lifting
Condition (\hyperref[def:yekutieli-lifting-condition]{Definition}).  In the
archimedean realisation, $\mathcal{B} = \mathcal{O}_M^{\times}$.  In
particular, the lift of the Frobenius--Hopf object from level $n$ to
level $n{+}1$ exists if and only if the obstruction class
$\kappa_n \in H^2(M,\, \mathcal{B})$
(\hyperref[def:quantum-characteristic-classes]{Definition}) vanishes.

---

By [Theorem 1.7, yekutieli2010], any central extension of gerbes
banded by an abelian sheaf on a ringed space satisfies the lifting
criterion: an object in the base gerbe admits a lift to the total gerbe
if and only if its obstruction class in $H^2$ vanishes.  The jet
prolongation construction
(\hyperref[prop:jet-prolongation]{Proposition}) realises the passage from
level $n$ to level $n{+}1$ as precisely such a central extension,
banded by $\mathcal{B}$
(\hyperref[def:quantum-characteristic-classes]{Definition}).  Hence the
Yekutieli Lifting Condition holds, and $\kappa_n = 0$ is equivalent
to the existence of the lift.
\end{proposition}
\begin{proof}
By [Theorem 1.7, yekutieli2010], any central extension of gerbes
banded by an abelian sheaf on a ringed space satisfies the lifting
criterion: an object in the base gerbe admits a lift to the total gerbe
if and only if its obstruction class in $H^2$ vanishes.  The jet
prolongation construction
(\hyperref[prop:jet-prolongation]{Proposition}) realises the passage from
level $n$ to level $n{+}1$ as precisely such a central extension,
banded by $\mathcal{B}$
(\hyperref[def:quantum-characteristic-classes]{Definition}).  Hence the
Yekutieli Lifting Condition holds, and $\kappa_n = 0$ is equivalent
to the existence of the lift.
\end{proof}

\begin{theorem}[Quantum universe with observability and Yekutieli is a QOU]
  \label{thm:qu-to-qou}
  \uses{def:quantum-universe, def:quantum-observable-universe, def:recursive-brane-tower, def:yekutieli-lifting-condition, def:quantum-characteristic-classes, def:projective-line-object, prop:jet-prolongation, prop:yekutieli-lifting, prop:fiber-adjoint-involution}
Let $(\mathcal{C}, A, \omega)$ be a quantum universe over a Noetherian
ring $R$ (\hyperref[def:quantum-universe]{Definition}); let $q$ denote the substrate parameter.  Assume observability
$\mathrm{Hom}(\mathbf{1},\, \mathrm{Ad}(A)) = 0$
(\hyperref[def:quantum-observable-universe]{Definition}).  Then the brane tower
$\{A_n\}_{n \geq 0}$ constructed via jet prolongation
(\hyperref[prop:jet-prolongation]{Proposition}) satisfies the Yekutieli lifting
condition (\hyperref[def:yekutieli-lifting-condition]{Definition},
\hyperref[prop:yekutieli-lifting]{Proposition}), and $(\mathcal{C}, A, \omega)$
is a quantum observable universe with the following structural properties:

\begin{enumerate}
  \item \emph{Descent terminates.}  The geometric descent tower has finite
\end{enumerate}
   length, bounded by $\mathrm{Krull\;dim}\,
   Z(\mathrm{End}_{\mathcal{C}}(A))$.  Observability ensures each step
   is a strict reduction (the $G_q$-action is faithful, so fixed-point
   reduction strictly decreases dimension).
\begin{enumerate}
  \item \emph{First obstruction equals the Euler class.}  The Yekutieli
\end{enumerate}
   obstruction $\kappa_0 \in H^2(\mathcal{S}, \mathcal{B})$ equals the
   Euler class of the Reeb bundle.  Under the trace map,
   $\kappa_0 \mapsto \chi(A) = \mathrm{tr}(\mathrm{id}_A) \in R$.
\begin{enumerate}
  \item \emph{Perpetual obstruction.}  If $\kappa_0$ is non-torsion in
\end{enumerate}
   $H^2(\mathcal{S}, \mathcal{B})$, then $\kappa_n \neq 0$ for all
   $n \geq 0$.  The class $\kappa_0$ and the jet corrections at
   level $n$ live in complementary subgroups of
   $H^2(\mathcal{S}, \mathcal{B})$, so the corrections cannot cancel
   $\kappa_0$.
\begin{enumerate}
  \item \emph{Branch point formula.}  The geometric descent from
\end{enumerate}
   $\mathbb{P}(V)$ (\hyperref[def:projective-line-object]{Definition}) to
   points has $B$ branch points with total ramification
   $R_{\mathrm{tot}} = 2d + \lvert\chi(A)\rvert$ (Riemann--Hurwitz),
   where $d$ is the degree of the covering.
\begin{enumerate}
  \item \emph{Terminal fiber data.}  At each branch point $p$, the QOU
\end{enumerate}
   reduces to fiber data $(A, \omega, q(p))$ with the involution
   $\sigma \colon \mathbb{P}(V) \hookrightarrow \mathrm{Ad}(A)$
   (\hyperref[prop:fiber-adjoint-involution]{Proposition}).

---

\hyperref[prop:yekutieli-lifting]{Proposition} establishes that the constructed
brane tower satisfies the Yekutieli lifting condition.  We prove each
structural property in turn.

\emph{Item (1): Descent terminates.}
The centre $Z(\mathrm{End}_{\mathcal{C}}(A))$ is a finitely generated
commutative $R$-algebra (since $R$ is Noetherian and $\omega$ is faithful),
hence has finite Krull dimension $d_0$.  Observability forces the $G_q$-action
on $A$ to be faithful: any $G_q$-invariant element of $\mathrm{Ad}(A)$ would
give a non-zero element of $\mathrm{Hom}(\mathbf{1}, \mathrm{Ad}(A))$,
contradicting the observability condition.  Each descent step (fixed-point
reduction under $G_q$) therefore strictly reduces the Krull dimension of the
centre.  After at most $d_0$ steps, the Krull dimension reaches $0$, and the
base is a discrete (Artinian) scheme.

\emph{Item (2): $\kappa_0$ equals the Euler class.}
The first Yekutieli obstruction $\kappa_0 \in H^2(\mathcal{S}, \mathcal{B})$
is the connecting homomorphism image of the bundle class $[A_0]$ under the
central extension $1 \to \mathcal{B} \to \mathfrak{G}_1 \to \mathfrak{G}_0
\to 1$ (\hyperref[prop:yekutieli-lifting]{Proposition}).  In the archimedean
realisation, this identifies with the Euler class $e$ of the Reeb
$\mathrm{U}(1)$-bundle via the Gysin sequence.  The trace map sends
$\kappa_0$ to $\chi(A) = \mathrm{tr}(\mathrm{id}_A)$.

\emph{Item (3): Perpetual obstruction.}
The class $\kappa_0$ lies in the image of $H^2(\mathcal{S}, \underline{R})
\to H^2(\mathcal{S}, \mathcal{B})$ (constant/topological classes), while jet
corrections at level $n$ lie in the image of $H^2(\mathcal{S},
\mathcal{O}_{\mathcal{S}}) \to H^2(\mathcal{S}, \mathcal{B})$ (analytic/algebraic
classes).  When $\kappa_0$ is non-torsion, these images are complementary
subgroups, so the corrections cannot cancel $\kappa_0$.  In the archimedean
realisation, the splitting follows from the exponential sequence; in the
algebraic setting, from the étale--analytic comparison.

\emph{Item (4): Riemann--Hurwitz.}
The branched covering $\pi \colon S \to \mathbb{P}(V)$ of degree $d$ has
total ramification $R_{\mathrm{tot}} = 2d + \lvert\chi(A)\rvert$ by
Riemann--Hurwitz: $\chi(S) = d \cdot \chi(\mathbb{P}(V)) - R_{\mathrm{tot}}$
with $\chi(\mathbb{P}(V)) = 2$.

\emph{Item (5): Terminal fiber data.}
At a branch point $p$, the Krull dimension of the centre is $0$, so the base
is a point.  The QOU data reduces to the fiber: the Frobenius--Hopf
algebra $A$ at parameter $q(p)$, with the involution
$\sigma \colon \mathbb{P}(V) \hookrightarrow \mathrm{Ad}(A)$ of
\hyperref[prop:fiber-adjoint-involution]{Proposition}.
\end{theorem}
\begin{proof}
\hyperref[prop:yekutieli-lifting]{Proposition} establishes that the constructed
brane tower satisfies the Yekutieli lifting condition.  We prove each
structural property in turn.

\emph{Item (1): Descent terminates.}
The centre $Z(\mathrm{End}_{\mathcal{C}}(A))$ is a finitely generated
commutative $R$-algebra (since $R$ is Noetherian and $\omega$ is faithful),
hence has finite Krull dimension $d_0$.  Observability forces the $G_q$-action
on $A$ to be faithful: any $G_q$-invariant element of $\mathrm{Ad}(A)$ would
give a non-zero element of $\mathrm{Hom}(\mathbf{1}, \mathrm{Ad}(A))$,
contradicting the observability condition.  Each descent step (fixed-point
reduction under $G_q$) therefore strictly reduces the Krull dimension of the
centre.  After at most $d_0$ steps, the Krull dimension reaches $0$, and the
base is a discrete (Artinian) scheme.

\emph{Item (2): $\kappa_0$ equals the Euler class.}
The first Yekutieli obstruction $\kappa_0 \in H^2(\mathcal{S}, \mathcal{B})$
is the connecting homomorphism image of the bundle class $[A_0]$ under the
central extension $1 \to \mathcal{B} \to \mathfrak{G}_1 \to \mathfrak{G}_0
\to 1$ (\hyperref[prop:yekutieli-lifting]{Proposition}).  In the archimedean
realisation, this identifies with the Euler class $e$ of the Reeb
$\mathrm{U}(1)$-bundle via the Gysin sequence.  The trace map sends
$\kappa_0$ to $\chi(A) = \mathrm{tr}(\mathrm{id}_A)$.

\emph{Item (3): Perpetual obstruction.}
The class $\kappa_0$ lies in the image of $H^2(\mathcal{S}, \underline{R})
\to H^2(\mathcal{S}, \mathcal{B})$ (constant/topological classes), while jet
corrections at level $n$ lie in the image of $H^2(\mathcal{S},
\mathcal{O}_{\mathcal{S}}) \to H^2(\mathcal{S}, \mathcal{B})$ (analytic/algebraic
classes).  When $\kappa_0$ is non-torsion, these images are complementary
subgroups, so the corrections cannot cancel $\kappa_0$.  In the archimedean
realisation, the splitting follows from the exponential sequence; in the
algebraic setting, from the étale--analytic comparison.

\emph{Item (4): Riemann--Hurwitz.}
The branched covering $\pi \colon S \to \mathbb{P}(V)$ of degree $d$ has
total ramification $R_{\mathrm{tot}} = 2d + \lvert\chi(A)\rvert$ by
Riemann--Hurwitz: $\chi(S) = d \cdot \chi(\mathbb{P}(V)) - R_{\mathrm{tot}}$
with $\chi(\mathbb{P}(V)) = 2$.

\emph{Item (5): Terminal fiber data.}
At a branch point $p$, the Krull dimension of the centre is $0$, so the base
is a point.  The QOU data reduces to the fiber: the Frobenius--Hopf
algebra $A$ at parameter $q(p)$, with the involution
$\sigma \colon \mathbb{P}(V) \hookrightarrow \mathrm{Ad}(A)$ of
\hyperref[prop:fiber-adjoint-involution]{Proposition}.
\end{proof}

\begin{remark}[Fermionic sectors as twisted sectors]
  \label{rem:fermionic-twisted-sectors}
  \uses{def:quantum-characteristic-classes, rem:algebraic-filter}
When the Yekutieli obstruction class $\kappa_n$ does not vanish, the
lift $E_n \to E_{n+1}$ is obstructed and the transition between brane
levels acquires a non-trivial twist.  In the algebraic filter
(\hyperref[rem:algebraic-filter]{Remark}), fermionic-type matter --- those
$q$-harmonic forms satisfying $[F_A,\, \alpha] \neq 0$ --- exists
precisely at the loci where $\kappa_n \neq 0$.  The non-vanishing
obstruction class creates "twisted" sectors in the gauge flow: the
holonomy around these sectors acquires a phase in
$H^2(M,\, \mathcal{O}_M^{\times})$ that prevents the vortex filament
from lifting smoothly, forcing the localisation of fermionic excitations
at the monodromy fixed points of the brane tower.
\end{remark}

\section{Betti Numbers and Harmonic Basis}
\label{sec:betti-harmonic}

\begin{remark}[Betti numbers of the brane tower]
  \label{rem:betti-numbers-brane}
  \uses{def:quantum-characteristic-classes}
Let $b_k(E_n) = \dim H^k(E_n, \mathbb{R})$ denote the $k$-th Betti number of
$E_n$.  Since $E_{n+1}$ is the bundle of orthonormal frames of $J^1 E_n$, the
Betti numbers of $E_{n+1}$ are determined by those of $E_n$ via the Leray--Serre
spectral sequence applied to the fibration
$E_{n+1} \to E_n$.  In particular, the quantum characteristic classes
(\hyperref[def:quantum-characteristic-classes]{Definition}) obstruct the collapse of
this spectral sequence and control the growth of $b_k(E_{n+1})$ in terms of
$b_k(E_n)$.
\end{remark}

\begin{remark}[$q$-Harmonic basis of the brane tower]
  \label{rem:q-harmonic-basis-brane}
  \uses{def:recursive-brane-tower, def:substrate-parameter, prop:jet-prolongation, rem:betti-numbers-brane}
For each level $n$, the $q$-harmonic $1$-forms $\alpha_{n,i}$ belong to the
space $\mathcal{H}^1(E_n,\, \mathrm{ad}\, E_n)$ of harmonic $1$-forms on
$E_n$ with values in the adjoint bundle $\mathrm{ad}\, E_n$.  These are
characterised by the self-duality condition relative to the instanton
connection $A_n$
(\hyperref[def:recursive-brane-tower]{Definition}):
\[
\alpha_{n,i} \in \mathcal{H}^1
  \iff
  d_{A_n}\,\alpha = 0
  \quad\text{and}\quad
  d_{A_n}^*\,\alpha = 0,
\]
where the Hodge-star operator is induced by the Hermitian metric that the
Levi form $L$ (\hyperref[def:substrate-parameter]{Definition}) defines on the
horizontal distribution.
The index $i$ runs over $1 \leq i \leq b_1(E_n)$
(\hyperref[rem:betti-numbers-brane]{Remark}).

At the base level $n = 0$, with Euler characteristic $\chi(M) = 3$, there are
three primary basis vectors $\{\,\phi_1,\, \phi_2,\, \phi_3\,\}$.  At
level $n = 1$ the basis lifts: the new basis includes the original modes
together with their $J^1$ prolongations
(\hyperref[prop:jet-prolongation]{Proposition}), which are the first-order
resolutions of the base flows.
\end{remark}

\begin{remark}[Weight-space decomposition of the harmonic basis]
  \label{rem:weight-space-decomposition}
  \uses{def:direct-limit-algebra, rem:j1e-local-states}
The adjoint bundle $\mathrm{ad}\, E_n$ has fibre $\mathfrak{su}(2)$.  Via the
weight-space decomposition of $\mathrm{SU}(2)$, the harmonic basis of
$\mathcal{H}^1(E_n,\, \mathrm{ad}\, E_n)$ splits into two families.

\emph{Longitudinal modes.}
These forms constitute the $\mathrm{U}(1)$ backbone of the decomposition.
They take values in the Cartan subalgebra of $\mathfrak{su}(2)$ and have the
geometric form $\omega^0 \otimes H$, where $H$ is the diagonal generator.
They represent steady-state laminar flows aligned with the $\mathrm{U}(1)$
subgroup.  If the bundle splits at a monodromy fixed point, these are the only
modes that survive; they correspond to the magnetic flux of the stable
vortices.

\emph{Ladder modes.}
These are the raising and lowering vortices of the decomposition.  They take
values in the off-diagonal generators $E_\pm$ of $\mathfrak{su}(2)$ and have
the geometric form $\omega^\pm \otimes E_\pm$.  They describe the entangling
flows that twist the $\mathbb{P}^1(\mathbb{C})$ fibre as one moves along the brane.
These modes are responsible for the non-commutativity in the
$\mathrm{SU}_q(2)$ algebra (\hyperref[def:direct-limit-algebra]{Definition}).

In the $1$-jet bundle $J^1 E_n$
(\hyperref[rem:j1e-local-states]{Remark}), each basis element $\alpha$ is a pair
consisting of a \emph{value} (the direction of the $\mathrm{SU}(2)$ flow in
the fibre) and a \emph{momentum} (the gradient of that flow relative to the
complex structure of the brane).
\end{remark}

\begin{remark}[The algebraic filter]
  \label{rem:algebraic-filter}
  \uses{rem:q-harmonic-basis-brane, def:substrate-parameter, rem:weight-space-decomposition}
Let $\alpha$ be a harmonic $1$-form in the basis of
\hyperref[rem:q-harmonic-basis-brane]{Remark}.  Its \emph{signature} is determined
by its commutator with the curvature $F_A$
(\hyperref[def:substrate-parameter]{Definition}).

\emph{Fermionic type (matter).}
$[F_A,\, \alpha] \neq 0$.  These forms are entangled with the curvature: they
see the $\mathrm{SO}(3)$ holonomy and the $-6$ index fixed points.  Their
geodesics are trapped or spiralling around the three fundamental vortex
filaments.

\emph{Bosonic type (radiation).}
$[F_A,\, \alpha] = 0$.  These forms commute with the curvature.  They are the
longitudinal modes (\hyperref[rem:weight-space-decomposition]{Remark}) that flow
through the bulk without being deflected by the monodromy fixed points.  Their
geodesics are straight relative to the gauge-fluid flow.
\end{remark}

\begin{remark}[Distributional connection at the limit]
  \label{rem:distributional-connection}
  \uses{rem:q-harmonic-basis-brane, rem:jet-shell}
At the limit $n \to \infty$ one no longer looks at a single connection $A$,
but at a \emph{distributional connection} represented by the series
\[
A_\infty
  \;=\;
  \sum_{n=0}^{\infty}\;\sum_{i=1}^{b_1(E_n)}
  c_{n,i}\,\alpha_{n,i}\,,
\]
where $c_{n,i} \in \mathbb{C}$ and $\{\alpha_{n,i}\}$ is the $q$-harmonic
basis of \hyperref[rem:q-harmonic-basis-brane]{Remark}.  This turns the gauge field
into a harmonic expansion across the infinite jet space $J^\infty E$
(\hyperref[rem:jet-shell]{Remark}).
\end{remark}

\section{Direct Limit Algebra}
\label{sec:direct-limit}

% ERROR: block "direct-limit-algebra" not found
\begin{remark}[Direct limit algebra and the Lusztig integral form]
  \label{rem:direct-limit-lusztig}
  \uses{def:direct-limit-algebra, def:substrate-parameter, def:quaternionic-spectral-resolvents}
The algebra $\mathcal{A}_\infty$
(\hyperref[def:direct-limit-algebra]{Definition}) is explicitly realised as the
Lusztig integral form of the quantised enveloping algebra
$\mathcal{U}_q(\mathfrak{su}(2))$ [lusztig1993].  At each level $n$
the fibre algebra $\mathrm{SU}_{q_n}(2)$ is obtained by specialising the
$\mathbb{Z}[q,q^{-1}]$-form
$\mathcal{U}^{\mathbb{Z}}_q(\mathfrak{su}(2))$
(\hyperref[def:substrate-parameter]{Definition}) along
$q \mapsto q_n \in k^{\times}$, and the direct limit inherits the divided
power basis $(E^{(m)}, F^{(m)}, K^{\pm 1})$ from the integral form.  This
ensures that the quaternionic jet embedding
(\hyperref[def:quaternionic-spectral-resolvents]{Definition}) remains robust
under the arithmetic descent to $M_0$, with the $q$-deformed scaling
function $\exp_q(r) - 1$ (equation \hyperref[eq:q-exp]{eq:q-exp}) regulating the
transition between the divided power basis of the space-like and time-like
resolvents.

The flatness of this specialisation---the fact that the divided-power
basis remains a basis after $q \mapsto q_n$---is guaranteed by the
Gröbner--Shirshov normal-form theorem applied to
$\mathrm{Dist}(G)$ [fred2005].  In the $p$-adic setting
(\hyperref[chap:p-adic]{Chapter}), the Witt vector lift
(\hyperref[def:witt-vector-lift]{Definition}) converts this algebraic
flatness into a geometric statement: the formal group $\widehat{G}$
over $\mathbb{Z}_p$ reduces modulo $p$ to the commutative
formal group $H$ over $\mathbb{F}_p$, and the divided-power basis
of $\mathrm{Dist}(\widehat{G})$ restricts to the PBW basis of
$\mathrm{Dist}(H)$.
\end{remark}

\begin{remark}[Jet Shell of the universe]
  \label{rem:jet-shell}
  \uses{def:direct-limit-algebra}
In the limit $n \to \infty$, the infinite jet bundle
$J^\infty E = \varprojlim\, J^k E$ becomes the \emph{jet shell} of the
universe.  All finite-level brane data are encoded in this pro-finite
structure, and the algebra $\mathcal{A}_\infty$
(\hyperref[def:direct-limit-algebra]{Definition}) acts naturally on sections of the
jet shell.
\end{remark}

\begin{remark}[Stack-like structure of the direct limit algebra]
  \label{rem:stack-like-structure}
  \uses{def:direct-limit-algebra, rem:jet-shell}
The direct limit algebra $\mathcal{A}_\infty$
(\hyperref[def:direct-limit-algebra]{Definition}) is a global "stack-like"
version of the local $\mathrm{SU}_q(2)$ deformations: at each point
$p \in M$ the fibre algebra $\mathcal{O}_{q(p)}(\mathrm{SU}(2))$
defines a local twisted deformation in the sense of
Yekutieli [yekutieli2008], and the direct limit assembles these
into a global object on the jet shell $J^\infty E$
(\hyperref[rem:jet-shell]{Remark}).
\end{remark}

\begin{proposition}[Global gluing via Yekutieli's twisted quantization]
  \label{prop:global-gluing}
  \uses{def:quantum-observable-universe, def:direct-limit-algebra, def:quantum-planck-constant, def:quantum-characteristic-classes, prop:vacuum-state}
Let $(\mathcal{C}, A, \omega)$ be a quantum observable universe
(\hyperref[def:quantum-observable-universe]{Definition}); let $q$ denote the substrate parameter.  The local
deformations of the Frobenius--Hopf object $A$ parameterised by the
substrate parameter $q$ induce a twisted associative deformation that
admits a unique global gluing via Yekutieli's twisted quantization
map [yekutieli2008, yekutieli2010].  In particular, the direct
limit algebra $\mathcal{A}_\infty$
(\hyperref[def:direct-limit-algebra]{Definition}) is a well-defined global
section of a sheaf of twisted deformation algebras on the underlying
site $\mathcal{S}$, and the vacuum state
(\hyperref[prop:vacuum-state]{Proposition}) is well-defined globally despite
local non-commutativity.  In the archimedean realisation, this
corresponds to the global gluing of the local algebras
$\mathcal{O}_{q(p)}(\mathrm{SU}(2))$ over the manifold $M$.

---

At each point $p \in M$ the fibre algebra
$\mathcal{O}_{q(p)}(\mathrm{SU}(2))$ is a pronilpotent deformation of
$\mathcal{O}(\mathrm{SU}(2))$ in the formal parameter
$\hbar_q(p) = 1 - q(p)^{-1}$
(\hyperref[def:quantum-planck-constant]{Definition}).
By [yekutieli2008], the coordinate bundle $\mathrm{Coor}_M$ of $M$
(the bundle of formal coordinate systems) carries a twisted deformation
datum: a pair consisting of a connection and a Maurer--Cartan element in
a pronilpotent DG Lie algebra [yekutieli2009].  The Yekutieli
obstruction classes $\kappa_n$
(\hyperref[def:quantum-characteristic-classes]{Definition}) control whether
these local twisted deformations glue across coordinate patches.

The vanishing condition
$\mathrm{Hom}(\mathbf{1}, \mathrm{Ad}(A)) = 0$
(\hyperref[def:quantum-observable-universe]{Definition}) ensures that the
gerbe $\mathfrak{G}_n$ of local lifts has trivial band at the level of
isomorphisms, so by [Theorem 4.5, yekutieli2010] the twisted
deformations at each fibre patch together into a global sheaf of
algebras on $M$.  The direct limit $\mathcal{A}_\infty =
\varinjlim\,\Gamma(J^1 E_n,\, \mathrm{SU}_{q_n}(2))$ is then the
global sections of this sheaf, and the vacuum state $\omega$
(\hyperref[prop:vacuum-state]{Proposition}) extends uniquely by the
$\mathrm{SU}(2)$-averaging argument of that proposition.
\end{proposition}
\begin{proof}
At each point $p \in M$ the fibre algebra
$\mathcal{O}_{q(p)}(\mathrm{SU}(2))$ is a pronilpotent deformation of
$\mathcal{O}(\mathrm{SU}(2))$ in the formal parameter
$\hbar_q(p) = 1 - q(p)^{-1}$
(\hyperref[def:quantum-planck-constant]{Definition}).
By [yekutieli2008], the coordinate bundle $\mathrm{Coor}_M$ of $M$
(the bundle of formal coordinate systems) carries a twisted deformation
datum: a pair consisting of a connection and a Maurer--Cartan element in
a pronilpotent DG Lie algebra [yekutieli2009].  The Yekutieli
obstruction classes $\kappa_n$
(\hyperref[def:quantum-characteristic-classes]{Definition}) control whether
these local twisted deformations glue across coordinate patches.

The vanishing condition
$\mathrm{Hom}(\mathbf{1}, \mathrm{Ad}(A)) = 0$
(\hyperref[def:quantum-observable-universe]{Definition}) ensures that the
gerbe $\mathfrak{G}_n$ of local lifts has trivial band at the level of
isomorphisms, so by [Theorem 4.5, yekutieli2010] the twisted
deformations at each fibre patch together into a global sheaf of
algebras on $M$.  The direct limit $\mathcal{A}_\infty =
\varinjlim\,\Gamma(J^1 E_n,\, \mathrm{SU}_{q_n}(2))$ is then the
global sections of this sheaf, and the vacuum state $\omega$
(\hyperref[prop:vacuum-state]{Proposition}) extends uniquely by the
$\mathrm{SU}(2)$-averaging argument of that proposition.
\end{proof}

\begin{proposition}[Existence and Uniqueness of the Vacuum State]
  \label{prop:vacuum-state}
  \uses{def:quantum-observable-universe, def:direct-limit-algebra, def:quantum-gauge-group, lem:descent}
Let $(\mathcal{C}, A, \omega)$ be a quantum observable universe
(\hyperref[def:quantum-observable-universe]{Definition}); let $q$ denote the substrate parameter.  Let
$\mathcal{A}_\infty$ be its direct limit algebra
(\hyperref[def:direct-limit-algebra]{Definition}).  Let $G_q =
\mathrm{Aut}^{\otimes}(\omega)$ be the quantum gauge group
(\hyperref[def:quantum-gauge-group]{Definition}), acting on
$\mathcal{A}_\infty$ by conjugation.  Then there exists a unique
$R$-linear functional
$\varphi \colon \mathcal{A}_\infty \to R$ that is $G_q$-invariant:
\[
\varphi(g \cdot a) = \varphi(a)
  \quad\text{for all } g \in G_q,\;
  a \in \mathcal{A}_\infty.
\]
In the archimedean realisation, $G_q = \mathrm{SU}_q(2)$ and
$\varphi$ is the unique $\mathrm{SU}(2)$-invariant state on the
$C^*$-algebraic direct limit.

---

The algebra $\mathcal{A}_\infty$ is the direct limit of the fibre algebras
$\Gamma(J^1 E_n,\, \mathrm{SU}_{q_n}(2))$
(\hyperref[def:direct-limit-algebra]{Definition}), each of which is a
$C^*$-algebra carrying a continuous $\mathrm{SU}(2)$ action by automorphisms.

\emph{Existence.}
At each finite level $n$ the algebra
$\Gamma(J^1 E_n,\, \mathrm{SU}_{q_n}(2))$ is a type-I $C^*$-algebra
(it is a section algebra of a locally trivial bundle with compact fibre).
By compactness of $\mathrm{SU}(2)$, averaging any state $\psi_n$ over the
Haar measure $\mu$ produces an $\mathrm{SU}(2)$-invariant state:
\[
\omega_n(a)
  \;=\;
  \int_{\mathrm{SU}(2)} \psi_n(g \cdot a)\; d\mu(g).
\]
The descent homomorphisms $\mathfrak{d}_n$ (\hyperref[lem:descent]{Lemma}) are
equivariant and surjective, so the invariant states $\{\omega_n\}$ form a
compatible system.  By the universal property of the direct limit, they
determine a state $\omega$ on $\mathcal{A}_\infty$ satisfying
$\omega(g \cdot a) = \omega(a)$.

\emph{Uniqueness.}
The vanishing condition
$\mathrm{Hom}(\mathbf{1}, \mathrm{Ad}(A)) = 0$ of the quantum observable
universe axioms implies that the trivial representation appears with
multiplicity one in each $\Gamma(J^1 E_n,\, \mathrm{SU}_{q_n}(2))$ (it
occurs only in the scalar summand $\mathbb{C} \cdot \mathbf{1}$).  Any
$\mathrm{SU}(2)$-invariant state must factor through the projection onto
this trivial isotypic component, which is one-dimensional.  Hence
$\omega$ is the unique such state.
\end{proposition}
\begin{proof}
The algebra $\mathcal{A}_\infty$ is the direct limit of the fibre algebras
$\Gamma(J^1 E_n,\, \mathrm{SU}_{q_n}(2))$
(\hyperref[def:direct-limit-algebra]{Definition}), each of which is a
$C^*$-algebra carrying a continuous $\mathrm{SU}(2)$ action by automorphisms.

\emph{Existence.}
At each finite level $n$ the algebra
$\Gamma(J^1 E_n,\, \mathrm{SU}_{q_n}(2))$ is a type-I $C^*$-algebra
(it is a section algebra of a locally trivial bundle with compact fibre).
By compactness of $\mathrm{SU}(2)$, averaging any state $\psi_n$ over the
Haar measure $\mu$ produces an $\mathrm{SU}(2)$-invariant state:
\[
\omega_n(a)
  \;=\;
  \int_{\mathrm{SU}(2)} \psi_n(g \cdot a)\; d\mu(g).
\]
The descent homomorphisms $\mathfrak{d}_n$ (\hyperref[lem:descent]{Lemma}) are
equivariant and surjective, so the invariant states $\{\omega_n\}$ form a
compatible system.  By the universal property of the direct limit, they
determine a state $\omega$ on $\mathcal{A}_\infty$ satisfying
$\omega(g \cdot a) = \omega(a)$.

\emph{Uniqueness.}
The vanishing condition
$\mathrm{Hom}(\mathbf{1}, \mathrm{Ad}(A)) = 0$ of the quantum observable
universe axioms implies that the trivial representation appears with
multiplicity one in each $\Gamma(J^1 E_n,\, \mathrm{SU}_{q_n}(2))$ (it
occurs only in the scalar summand $\mathbb{C} \cdot \mathbf{1}$).  Any
$\mathrm{SU}(2)$-invariant state must factor through the projection onto
this trivial isotypic component, which is one-dimensional.  Hence
$\omega$ is the unique such state.
\end{proof}

\section{Mass in the Brane Tower}
\label{sec:mass-brane}

\begin{remark}[Mass in the brane tower]
  \label{rem:mass-geometric-drag}
  \uses{def:recursive-brane-tower, def:substrate-parameter, prop:jet-prolongation}
In the recursive brane tower (\hyperref[def:recursive-brane-tower]{Definition}),
mass is not an intrinsic property of a point-particle but a measure of the
resistance encountered by a stable vortex as it propagates through the
$1$-jet bundle $J^1 E_n$ (\hyperref[prop:jet-prolongation]{Proposition}).  The
substrate parameter $q$ (\hyperref[def:substrate-parameter]{Definition}) varies from
fibre to fibre; translation of the vortex along the base therefore requires a
corresponding variation of the Frobenius--Hopf structure on each fibre, and
the mass quantifies this variation.
\end{remark}

\begin{remark}[Non-commutativity and geometric drag]
  \label{rem:mass-noncommutative-drag}
  \uses{def:quantum-universe, def:direct-limit-algebra, rem:mass-geometric-drag}
When the substrate parameter satisfies $q(p) \neq 1$ at every point of the recursive brane tower, the comultiplication on the Frobenius--Hopf object $A$
(\hyperref[def:quantum-universe]{Definition}) is non-cocommutative, so the
algebra $\mathcal{A}_\infty$ (\hyperref[def:direct-limit-algebra]{Definition})
does not reduce to a commutative direct limit.  Moving a vortex through
$J^1 E_n$ then requires a continuous rearrangement of the
$\mathrm{SU}_q(2)$ fibre data, which is the source of the inertial
resistance described in \hyperref[rem:mass-geometric-drag]{Remark}.
\end{remark}

\begin{proposition}[Stability of the Euler characteristic]
  \label{prop:euler-char-stability}
  \uses{def:quantum-observable-universe, def:recursive-brane-tower, def:quantum-characteristic-classes, prop:jet-prolongation}
Let $(\mathcal{C}, A, \omega)$ be a quantum observable universe
(\hyperref[def:quantum-observable-universe]{Definition}) and let
$\{A_n\}_{n \geq 0}$ be its recursive brane tower
(\hyperref[def:recursive-brane-tower]{Definition}).  The categorical Euler
characteristic $\chi(A_n) := \mathrm{tr}(\mathrm{id}_{A_n})$ (computed
via the rigidity trace of the identity endomorphism in $\mathcal{C}$)
is stable under jet prolongation: for each level $n$,
\[
\chi(A_{n+1}) \;=\; \chi(A_n).
\]
In particular, if the base level satisfies $\chi(A_0) = c \in k$, then
$\chi(A_n) = c$ for all $n \geq 0$.  In the archimedean realisation,
$\chi(A_n) = \chi(E_n)$ is the topological Euler characteristic of the
total space $E_n$.

---

By the construction of the recursive brane tower
(\hyperref[def:recursive-brane-tower]{Definition}), $E_{n+1}$ is the bundle of
orthonormal frames of the first jet bundle $J^1 E_n$
(\hyperref[prop:jet-prolongation]{Proposition}).  The Euler characteristic of a
fibre bundle $F \hookrightarrow E_{n+1} \to E_n$ satisfies the
multiplicativity formula
\[
\chi(E_{n+1}) \;=\; \chi(E_n) \cdot \chi(F),
\]
provided the fibre $F$ has the homotopy type of a finite CW-complex and the
base acts trivially on $H^*(F;\, \mathbb{Q})$.  In the brane tower the fibre
is $\mathrm{SU}(2) \cong S^3$, which has $\chi(S^3) = 0$ for the odd-dimensional
sphere.  However, the relevant bundle is not $\mathrm{SU}(2)$ itself but the
\emph{projectivised} frame bundle, whose typical fibre is
$\mathrm{SU}(2)/\mathrm{U}(1) \cong \mathbb{P}^1(\mathbb{C}) \cong S^2$ with
$\chi(S^2) = 2$.  Nonetheless, the Euler characteristic is preserved for the
following reason.

The Leray--Serre spectral sequence of
$S^2 \hookrightarrow E_{n+1} \to E_n$ has $E_2^{p,q} = H^p(E_n;\,
H^q(S^2;\, \mathbb{Z}))$.  The quantum characteristic classes $\kappa_n$
(\hyperref[def:quantum-characteristic-classes]{Definition}) are precisely the
obstructions to trivial monodromy of the fibre cohomology $H^*(S^2)$.  By
the vanishing condition of the quantum observable universe axioms---which
forces $\mathrm{Hom}(\mathbf{1}, \mathrm{Ad}(A)) = 0$---the monodromy
representation of $\pi_1(E_n)$ on $H^*(S^2;\, \mathbb{Z})$ is trivial (the
adjoint bundle has no invariant sections to twist the fibre cohomology).
Consequently the spectral sequence collapses at the $E_2$ page.  The
preservation of the Euler characteristic follows from the following
observation.  The base manifold $M_0$ is invariant under jet prolongation:
each level of the brane tower changes the bundle data but not the base
geometry.  Hence $\chi(M_0)$ is preserved at every level.
\end{proposition}
\begin{proof}
By the construction of the recursive brane tower
(\hyperref[def:recursive-brane-tower]{Definition}), $E_{n+1}$ is the bundle of
orthonormal frames of the first jet bundle $J^1 E_n$
(\hyperref[prop:jet-prolongation]{Proposition}).  The Euler characteristic of a
fibre bundle $F \hookrightarrow E_{n+1} \to E_n$ satisfies the
multiplicativity formula
\[
\chi(E_{n+1}) \;=\; \chi(E_n) \cdot \chi(F),
\]
provided the fibre $F$ has the homotopy type of a finite CW-complex and the
base acts trivially on $H^*(F;\, \mathbb{Q})$.  In the brane tower the fibre
is $\mathrm{SU}(2) \cong S^3$, which has $\chi(S^3) = 0$ for the odd-dimensional
sphere.  However, the relevant bundle is not $\mathrm{SU}(2)$ itself but the
\emph{projectivised} frame bundle, whose typical fibre is
$\mathrm{SU}(2)/\mathrm{U}(1) \cong \mathbb{P}^1(\mathbb{C}) \cong S^2$ with
$\chi(S^2) = 2$.  Nonetheless, the Euler characteristic is preserved for the
following reason.

The Leray--Serre spectral sequence of
$S^2 \hookrightarrow E_{n+1} \to E_n$ has $E_2^{p,q} = H^p(E_n;\,
H^q(S^2;\, \mathbb{Z}))$.  The quantum characteristic classes $\kappa_n$
(\hyperref[def:quantum-characteristic-classes]{Definition}) are precisely the
obstructions to trivial monodromy of the fibre cohomology $H^*(S^2)$.  By
the vanishing condition of the quantum observable universe axioms---which
forces $\mathrm{Hom}(\mathbf{1}, \mathrm{Ad}(A)) = 0$---the monodromy
representation of $\pi_1(E_n)$ on $H^*(S^2;\, \mathbb{Z})$ is trivial (the
adjoint bundle has no invariant sections to twist the fibre cohomology).
Consequently the spectral sequence collapses at the $E_2$ page.  The
preservation of the Euler characteristic follows from the following
observation.  The base manifold $M_0$ is invariant under jet prolongation:
each level of the brane tower changes the bundle data but not the base
geometry.  Hence $\chi(M_0)$ is preserved at every level.
\end{proof}

% ERROR: block "mass-identity" not found
\begin{corollary}[Archimedean realisation: nonabelian multiplicative integration]
  \label{cor:mass-identity-C}
  \uses{def:c-quantum-observable-universe, def:mass-identity}
In the archimedean quantum observable universe $(M, G, P, E)$
(\hyperref[def:c-quantum-observable-universe]{Example}) with instanton
connection $A$, Levi form $L$, and self-dual curvature $F^+$, the
categorical mass (\hyperref[def:mass-identity]{Definition}) is realised as
Yekutieli's nonabelian multiplicative integral [yekutieli2015].
Concretely, let $\gamma$ be a closed cycle encircling a vortex filament
and let $\Sigma$ be a surface with $\partial\Sigma = \gamma$.  Then:
\[
m(\gamma) \;=\;
  \left\lVert\,
  \mathrm{MI}\!\!\int_\Sigma
  \bigl(\nabla_A q,\; L \wedge F^+\bigr)
  \,\right\rVert,
\]
where $\mathrm{MI}\!\int_\Sigma$ denotes the $2$-dimensional twisted
nonabelian multiplicative integral, the integrand is the
connection--curvature pair $(\nabla_A q,\; L \wedge F^+)$ satisfying
the fake-flatness equation (\hyperref[eq:fake-flatness]{eq:fake-flatness}), and
$\lVert\cdot\rVert$ is the Killing-form norm on $\mathfrak{su}(2)$.
By the nonabelian Stokes theorem [yekutieli2015], this reduces to
the holonomy:
\[
m(\gamma) \;=\;
  \bigl\lVert\,
  \mathrm{Hol}_{\nabla_A}\!\bigl(\gamma;\, \mathfrak{G}\bigr)
  \,\bigr\rVert.
\]
\end{corollary}

\section{Yekutieli Categorical Mapping}
\label{sec:yekutieli-mapping}

\hyperref[tab:yekutieli-mapping]{Table} summarises the correspondence between the
core structures of the Quantum Observable Universe and the formal results
of Yekutieli on central extensions of gerbes, twisted deformation
quantization, and nonabelian multiplicative integration.

\begin{table}[ht]
\centering
\small
\begin{tabular}{p{4.2cm}p{4.8cm}p{4.5cm}}
\hline
\textbf{QOU concept} & \textbf{Yekutieli counterpart} & \textbf{Reference} \\
\hline
Recursive Brane Tower $\{E_n\}$
  (Def.~\ref{def:recursive-brane-tower})
  & Central extension of gerbes
    $\mathfrak{G}_{n+1} \to \mathfrak{G}_n$
  & \cite[Thm.~1.7]{yekutieli2010} \\[6pt]
Quantum Characteristic Classes $\kappa_n \in H^2(M, \mathcal{O}_M^{\times})$
  (Def.~\ref{def:quantum-characteristic-classes})
  & Obstruction classes for lifting objects in central extensions
  & \cite[\S3]{yekutieli2010} \\[6pt]
Fermionic sectors ($\kappa_n \neq 0$)
  (Rem.~\ref{rem:fermionic-twisted-sectors})
  & Twisted sectors where the lift is obstructed
  & \cite[Thm.~3.4]{yekutieli2010} \\[6pt]
Direct Limit Algebra $\mathcal{A}_\infty$
  (Def.~\ref{def:direct-limit-algebra})
  & Global sections of a sheaf of twisted deformation algebras
  & \cite{yekutieli2008} \\[6pt]
Global gluing of $\mathcal{A}_\infty$
  (Prop.~\ref{prop:global-gluing})
  & Yekutieli's twisted quantization map on algebraic varieties
  & \cite[Thm.~4.5]{yekutieli2010} \\[6pt]
Vacuum State $\omega$
  (Prop.~\ref{prop:vacuum-state})
  & Unique MC element in the pronilpotent DG Lie algebra
  & \cite{yekutieli2009} \\[6pt]
Mass Identity (Def.~\ref{def:mass-identity})
  & Nonabelian multiplicative integral; 2D nonabelian Stokes theorem
  & \cite[Ch.~5--6]{yekutieli2015} \\[6pt]
Substrate field $q(p)$
  (Def.~\ref{def:substrate-parameter})
  & Local twisted deformation parameter; MC element varying over $M$
  & \cite{yekutieli2008, yekutieli2009} \\[6pt]
$\mathrm{SU}_q(2)$ gauge flow
  (Def.~\ref{def:quantum-gauge-group})
  & Pronilpotent deformation of the coordinate algebra
  & \cite[\S2]{yekutieli2008} \\[6pt]
Descartes condition $\chi(E) = -6$
  (Def.~\ref{def:descartes-condition})
  & Topological constraint forcing $\kappa_n \neq 0$ (non-trivial gerbe)
  & \cite{yekutieli2010}; Ex.~\ref{ex:descartes-yekutieli} \\[6pt]
Cosimplicial crossed groupoid
  (Prop.~\ref{prop:cosimplicial-crossed-groupoid})
  & Cosimplicial object in crossed groupoids; \v{C}ech nerve of a gerbe
  & \cite[\S1--2]{yekutieli2010} \\[6pt]
Twisted $2$-cocycle $d(g_{ijk}) = a \cdot g \cdot a^{-1}$
  (Prop.~\ref{prop:descent-condition})
  & Twisted \v{C}ech cocycle in $H^2(M, \mathcal{O}_M^\times)$
  & \cite[Thm.~1.7]{yekutieli2010} \\
\hline
\end{tabular}
\caption{Categorical mapping between QOU structures and Yekutieli's results.}
\label{tab:yekutieli-mapping}
\end{table}

\section{Cosimplicial Crossed Groupoid Structure}
\label{ssec:cosimplicial-crossed-groupoid}

The recursive brane tower carries a higher-categorical structure that
organises the transitions between levels into a cosimplicial crossed
groupoid.  The key observation is that Yekutieli's framework for central
extensions of gerbes [yekutieli2010] furnishes exactly the triples
$(x, g, a)$ needed to define the face and degeneracy maps.

% ERROR: block "yekutieli-triple" not found
\begin{proposition}[Cosimplicial crossed groupoid]
  \label{prop:cosimplicial-crossed-groupoid}
  \uses{def:quantum-observable-universe, def:recursive-brane-tower, def:yekutieli-triple, prop:jet-prolongation, lem:descent}
Let $(\mathcal{C}, A, \omega)$ be a quantum observable universe and let
$\{A_n\}_{n \geq 0}$ be its recursive brane tower
(\hyperref[def:recursive-brane-tower]{Definition}); let $q$ denote the substrate parameter.  Equipped with the Yekutieli
triples $(x^{(n)}, g^{(n)}, a^{(n)})$ of
\hyperref[def:yekutieli-triple]{Definition}.  Then the assignment
\[
[n] \;\longmapsto\; \mathfrak{G}_n
\]
from the simplex category $\Delta$ to the $2$-category of gerbes on $M$
extends to a \emph{cosimplicial crossed groupoid}: a cosimplicial object
in the category of crossed groupoids, where

\begin{enumerate}
  \item the coface maps $\delta^i \colon \mathfrak{G}_{n} \to
\end{enumerate}
   \mathfrak{G}_{n+1}$ are given by the jet prolongation acting on
   objects (inserting the $i$-th degeneracy in the Cech nerve),
\begin{enumerate}
  \item the codegeneracy maps $\sigma^j \colon \mathfrak{G}_{n+1} \to
\end{enumerate}
   \mathfrak{G}_{n}$ are given by the descent homomorphism
   $\mathfrak{d}_{n+1}$ (collapsing the $j$-th face),
\begin{enumerate}
  \item the crossed structure is furnished by the conjugation action of
\end{enumerate}
   the $1$-morphisms $g_{ij}^{(n)}$ on the $2$-morphisms
   $a_{ij}^{(n)}$ via the adjoint representation of
   $\mathrm{SU}_{q_n}(2)$.

---

We verify the cosimplicial identities and the crossed-module axioms.

\emph{Step 1: Cosimplicial identities.}
The coface maps are composites of the jet prolongation with the
Cech face maps $\partial^i$ of the cover $\{U_i\}$.  The
codegeneracy maps are composites of the descent homomorphism with the
Cech degeneracy maps $s^j$.  The identities
$\delta^j \circ \delta^i = \delta^i \circ \delta^{j-1}$ for $i < j$
and $\sigma^j \circ \sigma^i = \sigma^i \circ \sigma^{j+1}$ for $i \leq j$
follow from the functoriality of jet prolongation and descent: the ascent
$\ell_n$ is a natural transformation from the identity functor on
$\mathfrak{G}_n$ to the restriction of $\mathfrak{G}_{n+1}$, and the
descent $\mathfrak{d}_{n+1}$ is its left adjoint.  The mixed identity
$\sigma^j \circ \delta^i$ reduces to three cases
($i < j$, $i = j$ or $i = j + 1$, $i > j + 1$), each verified using
the retraction property $\mathfrak{d}_{n+1} \circ \ell_n = \mathrm{id}$
(\hyperref[lem:descent]{Lemma}).

\emph{Step 2: Crossed-module structure.}
At each level $n$, the pair
$(\mathrm{SU}_{q_n}(2),\, \mathcal{O}_M^{\times})$ carries the structure
of a crossed module in the sense of [yekutieli2010]: the boundary
map $\partial \colon \mathcal{O}_M^{\times} \to \mathrm{SU}_{q_n}(2)$
sends a unit $a \in \mathcal{O}_M^{\times}$ to the central element
$a \cdot \mathbf{1} \in \mathrm{SU}_{q_n}(2)$ (the central extension),
and the action of $g \in \mathrm{SU}_{q_n}(2)$ on
$a \in \mathcal{O}_M^{\times}$ is conjugation:
$g \cdot a = g\, a\, g^{-1}$.  This action is well-defined because
$\mathcal{O}_M^{\times}$ is central in the gerbe band; the Peiffer
identity $\partial(a) \cdot b = a\, b\, a^{-1}$ is automatic for a
central extension.

\emph{Step 3: Compatibility.}
The coface and codegeneracy maps respect the crossed-module structure
because the jet prolongation preserves the Hopf-algebra multiplication
(\hyperref[prop:jet-prolongation]{Proposition}, Step 2) and the descent
homomorphism is a $*$-homomorphism (\hyperref[lem:descent]{Lemma}, Step 4).
Hence the conjugation action of $g_{ij}^{(n)}$ on $a_{ij}^{(n)}$ is
compatible with the simplicial maps, making the total assignment a
cosimplicial crossed groupoid.
\end{proposition}
\begin{proof}
We verify the cosimplicial identities and the crossed-module axioms.

\emph{Step 1: Cosimplicial identities.}
The coface maps are composites of the jet prolongation with the
Cech face maps $\partial^i$ of the cover $\{U_i\}$.  The
codegeneracy maps are composites of the descent homomorphism with the
Cech degeneracy maps $s^j$.  The identities
$\delta^j \circ \delta^i = \delta^i \circ \delta^{j-1}$ for $i < j$
and $\sigma^j \circ \sigma^i = \sigma^i \circ \sigma^{j+1}$ for $i \leq j$
follow from the functoriality of jet prolongation and descent: the ascent
$\ell_n$ is a natural transformation from the identity functor on
$\mathfrak{G}_n$ to the restriction of $\mathfrak{G}_{n+1}$, and the
descent $\mathfrak{d}_{n+1}$ is its left adjoint.  The mixed identity
$\sigma^j \circ \delta^i$ reduces to three cases
($i < j$, $i = j$ or $i = j + 1$, $i > j + 1$), each verified using
the retraction property $\mathfrak{d}_{n+1} \circ \ell_n = \mathrm{id}$
(\hyperref[lem:descent]{Lemma}).

\emph{Step 2: Crossed-module structure.}
At each level $n$, the pair
$(\mathrm{SU}_{q_n}(2),\, \mathcal{O}_M^{\times})$ carries the structure
of a crossed module in the sense of [yekutieli2010]: the boundary
map $\partial \colon \mathcal{O}_M^{\times} \to \mathrm{SU}_{q_n}(2)$
sends a unit $a \in \mathcal{O}_M^{\times}$ to the central element
$a \cdot \mathbf{1} \in \mathrm{SU}_{q_n}(2)$ (the central extension),
and the action of $g \in \mathrm{SU}_{q_n}(2)$ on
$a \in \mathcal{O}_M^{\times}$ is conjugation:
$g \cdot a = g\, a\, g^{-1}$.  This action is well-defined because
$\mathcal{O}_M^{\times}$ is central in the gerbe band; the Peiffer
identity $\partial(a) \cdot b = a\, b\, a^{-1}$ is automatic for a
central extension.

\emph{Step 3: Compatibility.}
The coface and codegeneracy maps respect the crossed-module structure
because the jet prolongation preserves the Hopf-algebra multiplication
(\hyperref[prop:jet-prolongation]{Proposition}, Step 2) and the descent
homomorphism is a $*$-homomorphism (\hyperref[lem:descent]{Lemma}, Step 4).
Hence the conjugation action of $g_{ij}^{(n)}$ on $a_{ij}^{(n)}$ is
compatible with the simplicial maps, making the total assignment a
cosimplicial crossed groupoid.
\end{proof}

\begin{remark}[Jet prolongation as formal descent of information]
  \label{rem:jet-as-descent}
  \uses{prop:cosimplicial-crossed-groupoid, prop:jet-prolongation, def:descent-involution}
In the cosimplicial crossed groupoid of
\hyperref[prop:cosimplicial-crossed-groupoid]{Proposition}, the coface map
$\delta^0 \colon \mathfrak{G}_n \to \mathfrak{G}_{n+1}$ applied to a
local section $x_i^{(n)}$ produces the jet-prolonged section
$x_i^{(n+1)} = J^1(x_i^{(n)})$.  This is precisely the content of
\hyperref[prop:jet-prolongation]{Proposition}: the jet prolongation is a
\emph{formal descent of information} from the higher categorical level
$\mathfrak{G}_{n+1}$ to the lower level $\mathfrak{G}_n$, packaged as a
coface map in the cosimplicial structure.  The fact that the codegeneracy
$\sigma^0 = \mathfrak{d}_{n+1}$ inverts this---in the sense that
$\sigma^0 \circ \delta^0 = \mathrm{id}$---is the retraction property that
underpins the descent involution (\hyperref[def:descent-involution]{Definition}).
\end{remark}

\begin{proposition}[Descent condition --- twisted $2$-cocycle]
  \label{prop:descent-condition}
  \uses{def:quantum-observable-universe, def:recursive-brane-tower, def:yekutieli-triple, def:quantum-characteristic-classes, rem:distributional-connection, prop:yekutieli-lifting}
Let $(\mathcal{C}, A, \omega)$ be a quantum observable universe and let
$\{A_n\}_{n \geq 0}$ be its recursive brane tower; let $q$ denote the substrate parameter.  With Yekutieli
triples $(x^{(n)}, g^{(n)}, a^{(n)})$
(\hyperref[def:yekutieli-triple]{Definition}) and let
$A_\infty = \sum_{n,i} c_{n,i}\, \alpha_{n,i}$
(\hyperref[rem:distributional-connection]{Remark}) be the distributional
connection.  Then the global consistency of $A_\infty$ is equivalent to
the satisfaction of the \emph{twisted $2$-cocycle condition}: for every
triple intersection $U_i \cap U_j \cap U_k$,

\begin{equation}\label{eq:twisted-2-cocycle}
  d\!\left(g_{ijk}^{(n)}\right)
  \;=\;
  a_{ij}^{(n)} \cdot g_{ijk}^{(n)} \cdot \left(a_{jk}^{(n)}\right)^{-1},
\end{equation}

where
$g_{ijk}^{(n)} = g_{ij}^{(n)} \cdot g_{jk}^{(n)} \cdot (g_{ik}^{(n)})^{-1}$
is the Cech $1$-cocycle deviation at level $n$,
$d(g_{ijk}^{(n)})$ denotes the Cech coboundary
(the failure of the $1$-cocycle condition), and
$a_{ij}^{(n)}, a_{jk}^{(n)} \in \mathcal{O}_M^{\times}$ are the
$2$-morphisms (quantum deformation feedback from $\mathrm{SU}_{q_n}(2)$).

Equivalently, the obstruction class $\kappa_n$
(\hyperref[def:quantum-characteristic-classes]{Definition}) is the cohomology
class of the twisted $2$-cocycle
$[a_{ij}^{(n)} \cdot (\cdot) \cdot (a_{jk}^{(n)})^{-1}]$ in
$H^2(M,\, \mathcal{O}_M^{\times})$.

---

Let $\{U_i\}$ be a good cover trivialising the tower.  On
$U_i \cap U_j$ the gauge transition function $g_{ij}^{(n)}$ patches
the connection $A_n$:
\[
A_n|_{U_j} \;=\; g_{ij}^{(n)} \cdot A_n|_{U_i} \cdot
  (g_{ij}^{(n)})^{-1} + (g_{ij}^{(n)})^{-1}\, dg_{ij}^{(n)}.
\]
On a triple overlap $U_i \cap U_j \cap U_k$, the ordinary cocycle
condition $g_{ij}^{(n)} \cdot g_{jk}^{(n)} = g_{ik}^{(n)}$ holds if and
only if the connection $A_n$ is globally well-defined at level $n$ with
no twist.  The passage from level $n$ to level $n{+}1$ via the central
extension $\mathfrak{G}_{n+1} \to \mathfrak{G}_n$ introduces the
$2$-morphism twist: the transition functions at level $n{+}1$ are related
to those at level $n$ by
\[
g_{ij}^{(n+1)} \;=\; a_{ij}^{(n)} \cdot g_{ij}^{(n)},
\]
where $a_{ij}^{(n)} \in \mathcal{O}_M^{\times}$ is the local lift datum
(\hyperref[def:yekutieli-triple]{Definition}).  Computing the cocycle deviation
at level $n{+}1$:

\begin{align*}
  g_{ijk}^{(n+1)}
  &= g_{ij}^{(n+1)} \cdot g_{jk}^{(n+1)} \cdot (g_{ik}^{(n+1)})^{-1} \\
  &= (a_{ij}^{(n)} \cdot g_{ij}^{(n)}) \cdot (a_{jk}^{(n)} \cdot g_{jk}^{(n)})
     \cdot (a_{ik}^{(n)} \cdot g_{ik}^{(n)})^{-1}.
\end{align*}

Since $a_{ij}^{(n)}$ is central (it lies in the band
$\mathcal{O}_M^{\times}$ of the gerbe), the gauge factors commute past
the $a$-factors, giving
\[
g_{ijk}^{(n+1)}
  \;=\;
  a_{ij}^{(n)} \cdot a_{jk}^{(n)} \cdot (a_{ik}^{(n)})^{-1}
  \cdot g_{ijk}^{(n)}.
\]
The distributional connection $A_\infty$ is globally consistent if and
only if $g_{ijk}^{(n+1)} = 0$ for all $n$ in the Cech cohomology,
which requires
\[
d(g_{ijk}^{(n)})
  \;=\;
  a_{ij}^{(n)} \cdot g_{ijk}^{(n)} \cdot (a_{jk}^{(n)})^{-1},
\]
where the left-hand side is the Cech coboundary
$\delta^2(g^{(n)})_{ijk}$ and the right-hand side is the twisted
conjugation by the $2$-morphisms.  The class of this twisted $2$-cocycle
in $H^2(M,\, \mathcal{O}_M^{\times})$ is precisely the obstruction
class $\kappa_n$ of
\hyperref[def:quantum-characteristic-classes]{Definition}: it vanishes if and
only if the cocycle deviation at level $n$ can be absorbed by the
$2$-morphism data, i.e. the lift from $\mathfrak{G}_n$ to
$\mathfrak{G}_{n+1}$ is unobstructed.

By [Theorem 1.7, yekutieli2010], this is equivalent to the
vanishing of the Yekutieli obstruction class, closing the circle with
the lifting criterion (\hyperref[prop:yekutieli-lifting]{Proposition}).
\end{proposition}
\begin{proof}
Let $\{U_i\}$ be a good cover trivialising the tower.  On
$U_i \cap U_j$ the gauge transition function $g_{ij}^{(n)}$ patches
the connection $A_n$:
\[
A_n|_{U_j} \;=\; g_{ij}^{(n)} \cdot A_n|_{U_i} \cdot
  (g_{ij}^{(n)})^{-1} + (g_{ij}^{(n)})^{-1}\, dg_{ij}^{(n)}.
\]
On a triple overlap $U_i \cap U_j \cap U_k$, the ordinary cocycle
condition $g_{ij}^{(n)} \cdot g_{jk}^{(n)} = g_{ik}^{(n)}$ holds if and
only if the connection $A_n$ is globally well-defined at level $n$ with
no twist.  The passage from level $n$ to level $n{+}1$ via the central
extension $\mathfrak{G}_{n+1} \to \mathfrak{G}_n$ introduces the
$2$-morphism twist: the transition functions at level $n{+}1$ are related
to those at level $n$ by
\[
g_{ij}^{(n+1)} \;=\; a_{ij}^{(n)} \cdot g_{ij}^{(n)},
\]
where $a_{ij}^{(n)} \in \mathcal{O}_M^{\times}$ is the local lift datum
(\hyperref[def:yekutieli-triple]{Definition}).  Computing the cocycle deviation
at level $n{+}1$:

\begin{align*}
  g_{ijk}^{(n+1)}
  &= g_{ij}^{(n+1)} \cdot g_{jk}^{(n+1)} \cdot (g_{ik}^{(n+1)})^{-1} \\
  &= (a_{ij}^{(n)} \cdot g_{ij}^{(n)}) \cdot (a_{jk}^{(n)} \cdot g_{jk}^{(n)})
     \cdot (a_{ik}^{(n)} \cdot g_{ik}^{(n)})^{-1}.
\end{align*}

Since $a_{ij}^{(n)}$ is central (it lies in the band
$\mathcal{O}_M^{\times}$ of the gerbe), the gauge factors commute past
the $a$-factors, giving
\[
g_{ijk}^{(n+1)}
  \;=\;
  a_{ij}^{(n)} \cdot a_{jk}^{(n)} \cdot (a_{ik}^{(n)})^{-1}
  \cdot g_{ijk}^{(n)}.
\]
The distributional connection $A_\infty$ is globally consistent if and
only if $g_{ijk}^{(n+1)} = 0$ for all $n$ in the Cech cohomology,
which requires
\[
d(g_{ijk}^{(n)})
  \;=\;
  a_{ij}^{(n)} \cdot g_{ijk}^{(n)} \cdot (a_{jk}^{(n)})^{-1},
\]
where the left-hand side is the Cech coboundary
$\delta^2(g^{(n)})_{ijk}$ and the right-hand side is the twisted
conjugation by the $2$-morphisms.  The class of this twisted $2$-cocycle
in $H^2(M,\, \mathcal{O}_M^{\times})$ is precisely the obstruction
class $\kappa_n$ of
\hyperref[def:quantum-characteristic-classes]{Definition}: it vanishes if and
only if the cocycle deviation at level $n$ can be absorbed by the
$2$-morphism data, i.e. the lift from $\mathfrak{G}_n$ to
$\mathfrak{G}_{n+1}$ is unobstructed.

By [Theorem 1.7, yekutieli2010], this is equivalent to the
vanishing of the Yekutieli obstruction class, closing the circle with
the lifting criterion (\hyperref[prop:yekutieli-lifting]{Proposition}).
\end{proof}

\begin{proposition}[Weak equivalence invariance of physical states]
  \label{prop:weak-equiv-invariance}
  \uses{def:quantum-observable-universe, def:recursive-brane-tower, def:direct-limit-algebra, def:descent-involution, prop:cosimplicial-crossed-groupoid, prop:vacuum-state, prop:descent-eigenspaces, lem:descent}
Let $(\mathcal{C}, A, \omega)$ be a quantum observable universe and let
$\{A_n\}_{n \geq 0}$ be its recursive brane tower with cosimplicial
crossed groupoid structure
(\hyperref[prop:cosimplicial-crossed-groupoid]{Proposition}); let $q$ denote the substrate parameter.  Let
$\omega \colon \mathcal{A}_\infty \to \mathbb{C}$ be the vacuum state
(\hyperref[prop:vacuum-state]{Proposition}).  A \emph{weak equivalence} between
brane levels $n$ and $m$ is a zigzag of coface and codegeneracy maps
\[
\mathfrak{G}_n
  \xleftarrow{\;\sigma\;}
  \mathfrak{G}_{k_1}
  \xrightarrow{\;\delta\;}
  \cdots
  \xrightarrow{\;\delta\;}
  \mathfrak{G}_m
\]
in the cosimplicial crossed groupoid.  Then for every observable
$a \in \mathcal{A}_\infty$, the expectation value is invariant under
weak equivalence:
\[
\omega\!\bigl(\iota_n(a|_n)\bigr)
  \;=\;
  \omega\!\bigl(\iota_m(a|_m)\bigr),
\]
where $\iota_n \colon \Gamma(J^1 E_n,\, \mathrm{SU}_{q_n}(2))
\hookrightarrow \mathcal{A}_\infty$ is the canonical embedding and
$a|_n$ denotes the restriction of $a$ to level $n$.  In particular,
physical states (the GNS vectors of $\omega$) are invariant under
fluctuations of the substrate parameter $q(p)$ across brane levels.

---

By [Theorem 1.1, yekutieli2010], a weak equivalence of central
extensions of gerbes induces an isomorphism on the category of global
objects up to isomorphism.  In the brane tower, the coface maps (jet
prolongation) and codegeneracy maps (descent) compose to give the descent
involution $\tau_n$ (\hyperref[def:descent-involution]{Definition}), which is
idempotent.  The image $\mathcal{A}_{n+1}^{\mathrm{desc}}$
(\hyperref[prop:descent-eigenspaces]{Proposition}) is therefore isomorphic, as
an algebra, to the level-$n$ section algebra via the ascent embedding
$\ell_n$.

For the vacuum state, the compatibility
$\omega(\iota_n(a|_n)) = \omega_n(a|_n)$ holds by the universal property
of the direct limit (\hyperref[def:direct-limit-algebra]{Definition}).  The
Haar-averaging construction in the proof of
\hyperref[prop:vacuum-state]{Proposition} produces states $\omega_n$ that are
compatible with the transition maps:
$\omega_{n+1} \circ \iota_{n,n+1} = \omega_n$, where
$\iota_{n,n+1}$ is the transition map.  A zigzag of coface and
codegeneracy maps between levels $n$ and $m$ factors as a sequence of
such transition maps and their adjoints (the descent homomorphisms),
each of which preserves the Haar state by equivariance
(\hyperref[lem:descent]{Lemma}, Step 5).

Therefore the expectation value
$\omega(\iota_n(a|_n)) = \omega_n(a|_n) = \omega_m(a|_m) =
\omega(\iota_m(a|_m))$
is invariant under weak equivalence.

For the GNS representation: the Hilbert space
$\mathcal{H}_\omega = \overline{\mathcal{A}_\infty / \ker(\omega)}$ with
inner product $\langle [a], [b] \rangle = \omega(a^* b)$ is unchanged by
the weak equivalence because $\omega$ is invariant.  The cyclic vector
$\Omega = [1]$ generates the same representation regardless of the brane
level at which it is evaluated, so the physical states (vectors in
$\mathcal{H}_\omega$) are insensitive to the substrate parameter
fluctuations across the tower.
\end{proposition}
\begin{proof}
By [Theorem 1.1, yekutieli2010], a weak equivalence of central
extensions of gerbes induces an isomorphism on the category of global
objects up to isomorphism.  In the brane tower, the coface maps (jet
prolongation) and codegeneracy maps (descent) compose to give the descent
involution $\tau_n$ (\hyperref[def:descent-involution]{Definition}), which is
idempotent.  The image $\mathcal{A}_{n+1}^{\mathrm{desc}}$
(\hyperref[prop:descent-eigenspaces]{Proposition}) is therefore isomorphic, as
an algebra, to the level-$n$ section algebra via the ascent embedding
$\ell_n$.

For the vacuum state, the compatibility
$\omega(\iota_n(a|_n)) = \omega_n(a|_n)$ holds by the universal property
of the direct limit (\hyperref[def:direct-limit-algebra]{Definition}).  The
Haar-averaging construction in the proof of
\hyperref[prop:vacuum-state]{Proposition} produces states $\omega_n$ that are
compatible with the transition maps:
$\omega_{n+1} \circ \iota_{n,n+1} = \omega_n$, where
$\iota_{n,n+1}$ is the transition map.  A zigzag of coface and
codegeneracy maps between levels $n$ and $m$ factors as a sequence of
such transition maps and their adjoints (the descent homomorphisms),
each of which preserves the Haar state by equivariance
(\hyperref[lem:descent]{Lemma}, Step 5).

Therefore the expectation value
$\omega(\iota_n(a|_n)) = \omega_n(a|_n) = \omega_m(a|_m) =
\omega(\iota_m(a|_m))$
is invariant under weak equivalence.

For the GNS representation: the Hilbert space
$\mathcal{H}_\omega = \overline{\mathcal{A}_\infty / \ker(\omega)}$ with
inner product $\langle [a], [b] \rangle = \omega(a^* b)$ is unchanged by
the weak equivalence because $\omega$ is invariant.  The cyclic vector
$\Omega = [1]$ generates the same representation regardless of the brane
level at which it is evaluated, so the physical states (vectors in
$\mathcal{H}_\omega$) are insensitive to the substrate parameter
fluctuations across the tower.
\end{proof}

\section{Descent}
\label{ssec:descent}



% ERROR: block "descent" not found
\begin{lemma}[Descent preserves the quantum observable universe]
  \label{lem:descent}
  \uses{def:quantum-observable-universe, def:recursive-brane-tower, prop:jet-prolongation}
Let $(\mathcal{C}, A, \omega)$ be a quantum observable universe and let
$\{A_n\}_{n \geq 0}$ be the recursive brane tower
(\hyperref[def:recursive-brane-tower]{Definition}), with $n \geq 1$; let $q$ denote the substrate parameter.  The
restriction morphism $r_n \colon A_n \to A_{n-1}$ yields a Frobenius--Hopf
object $A_{n-1}$ such that $(\mathcal{C}, A_{n-1}, \omega)$ is again a
quantum observable universe.  Moreover, the induced descent homomorphism
\[
\mathfrak{d}_n \;\colon\;
  \mathrm{End}_{\mathcal{C}}(A_n)
  \;\longrightarrow\;
  \mathrm{End}_{\mathcal{C}}(A_{n-1})
\]
is a surjective algebra homomorphism that intertwines the $G_q$-actions at
levels $n$ and $n{-}1$.  In the archimedean realisation, this recovers
the surjective $*$-homomorphism
$\Gamma(J^1 E_n,\, \mathrm{SU}_{q_n}(2)) \to
\Gamma(J^1 E_{n-1},\, \mathrm{SU}_{q_{n-1}}(2))$.

---

The proof mirrors the ascent (\hyperref[prop:jet-prolongation]{Proposition}) in
reverse, working inside the rigid monoidal category $\mathcal{C}$.

\emph{Step 1: Frobenius restriction.}
The jet object $J^1(A) = A \otimes \Omega^1(A)$ at level $n$ contains $A$
as a direct summand via the unit section
$\eta \otimes \mathrm{id}\colon \Omega^1(A) \hookrightarrow A \otimes
\Omega^1(A)$.  Composing with the counit $\varepsilon\colon A \to \mathbf{1}$
on the first tensor factor gives the restriction map
$r_n\colon J^1(A) \to A$.  Because $A$ is Frobenius, $r_n$ is a coalgebra
morphism: the comultiplication $\Delta$ commutes with restriction by the
Frobenius relation, and the counit restricts by construction.  Hence the
restricted object inherits a Frobenius--Hopf structure.

\emph{Step 2: Connection restriction.}
The connection $A_n$ on $E_n$ is determined by the comultiplication
$\Delta'$ on $J^1(A)$ (\hyperref[prop:jet-prolongation]{Proposition}, Step 2).
Restricting $\Delta'$ along $r_n$ recovers the original comultiplication
$\Delta$ on $A$, which by the construction of the brane tower corresponds to
the connection $A_{n-1}$ on $E_{n-1}$.  In particular, $A_{n-1}$ satisfies
the Hermitian Yang--Mills equations on $E_{n-1}$ because it did so at level
$n{-}1$ before the ascent.

\emph{Step 3: Non-degeneracy and vanishing.}
Non-degeneracy of the Frobenius form on the restricted object follows from
the direct-summand inclusion $A \hookrightarrow J^1(A)$: the restriction of a
non-degenerate pairing to a direct summand remains non-degenerate.  The
vanishing condition $\mathrm{Hom}(\mathbf{1}, \mathrm{Ad}(A)) = 0$ is
unchanged since we recover the original object $A$.

\emph{Step 4: The descent homomorphism.}
Applying the fibre functor $\omega$ to $r_n$ yields
$\mathfrak{d}_n\colon
\Gamma(J^1 E_n,\, \mathrm{SU}_{q_n}(2)) \to
\Gamma(J^1 E_{n-1},\, \mathrm{SU}_{q_{n-1}}(2))$.
This is a $*$-homomorphism because $r_n$ preserves multiplication and
the antipode.  Surjectivity follows from the fact that the ascent map
$\ell_n$ of \hyperref[prop:jet-prolongation]{Proposition} is injective (it
embeds $A$ into $J^1(A)$ as a direct summand), so its categorical adjoint
$r_n$ is surjective.

\emph{Step 5: Equivariance.}
The $\mathrm{SU}(2)$ action on sections at level $n$ is induced by the
adjoint representation of the Hopf algebra.  Since $r_n$ is a Hopf algebra
morphism, it intertwines the adjoint actions at levels $n$ and $n{-}1$.
\end{lemma}
\begin{proof}
The proof mirrors the ascent (\hyperref[prop:jet-prolongation]{Proposition}) in
reverse, working inside the rigid monoidal category $\mathcal{C}$.

\emph{Step 1: Frobenius restriction.}
The jet object $J^1(A) = A \otimes \Omega^1(A)$ at level $n$ contains $A$
as a direct summand via the unit section
$\eta \otimes \mathrm{id}\colon \Omega^1(A) \hookrightarrow A \otimes
\Omega^1(A)$.  Composing with the counit $\varepsilon\colon A \to \mathbf{1}$
on the first tensor factor gives the restriction map
$r_n\colon J^1(A) \to A$.  Because $A$ is Frobenius, $r_n$ is a coalgebra
morphism: the comultiplication $\Delta$ commutes with restriction by the
Frobenius relation, and the counit restricts by construction.  Hence the
restricted object inherits a Frobenius--Hopf structure.

\emph{Step 2: Connection restriction.}
The connection $A_n$ on $E_n$ is determined by the comultiplication
$\Delta'$ on $J^1(A)$ (\hyperref[prop:jet-prolongation]{Proposition}, Step 2).
Restricting $\Delta'$ along $r_n$ recovers the original comultiplication
$\Delta$ on $A$, which by the construction of the brane tower corresponds to
the connection $A_{n-1}$ on $E_{n-1}$.  In particular, $A_{n-1}$ satisfies
the Hermitian Yang--Mills equations on $E_{n-1}$ because it did so at level
$n{-}1$ before the ascent.

\emph{Step 3: Non-degeneracy and vanishing.}
Non-degeneracy of the Frobenius form on the restricted object follows from
the direct-summand inclusion $A \hookrightarrow J^1(A)$: the restriction of a
non-degenerate pairing to a direct summand remains non-degenerate.  The
vanishing condition $\mathrm{Hom}(\mathbf{1}, \mathrm{Ad}(A)) = 0$ is
unchanged since we recover the original object $A$.

\emph{Step 4: The descent homomorphism.}
Applying the fibre functor $\omega$ to $r_n$ yields
$\mathfrak{d}_n\colon
\Gamma(J^1 E_n,\, \mathrm{SU}_{q_n}(2)) \to
\Gamma(J^1 E_{n-1},\, \mathrm{SU}_{q_{n-1}}(2))$.
This is a $*$-homomorphism because $r_n$ preserves multiplication and
the antipode.  Surjectivity follows from the fact that the ascent map
$\ell_n$ of \hyperref[prop:jet-prolongation]{Proposition} is injective (it
embeds $A$ into $J^1(A)$ as a direct summand), so its categorical adjoint
$r_n$ is surjective.

\emph{Step 5: Equivariance.}
The $\mathrm{SU}(2)$ action on sections at level $n$ is induced by the
adjoint representation of the Hopf algebra.  Since $r_n$ is a Hopf algebra
morphism, it intertwines the adjoint actions at levels $n$ and $n{-}1$.
\end{proof}

% ERROR: block "descent-involution" not found
\begin{proposition}[Eigenspaces of the descent involution]
  \label{prop:descent-eigenspaces}
  \uses{def:descent-involution, rem:algebraic-filter, rem:weight-space-decomposition, lem:descent}
The descent involution $\tau_n$ (\hyperref[def:descent-involution]{Definition})
decomposes the level-$(n{+}1)$ endomorphism algebra into a direct sum
\[
\mathrm{End}_{\mathcal{C}}(A_{n+1})
  \;=\;
  \mathcal{A}_{n+1}^{\mathrm{desc}} \;\oplus\; \mathcal{A}_{n+1}^{\mathrm{trans}},
\]
where $\mathcal{A}_{n+1}^{\mathrm{desc}} = \mathrm{im}\,\tau_n$ is the
descended subalgebra (endomorphisms that factor through level $n$) and
$\mathcal{A}_{n+1}^{\mathrm{trans}} = \ker\,\tau_n$ is the *transient
complement* (genuinely new endomorphisms at level $n{+}1$).  The algebraic filter
(\hyperref[rem:algebraic-filter]{Remark}) refines this decomposition:

\begin{enumerate}
  \item The bosonic (longitudinal) modes lifted from level $n$ lie entirely in
\end{enumerate}
   $\mathcal{A}_{n+1}^{\mathrm{desc}}$: they commute with the curvature
   and survive every round-trip through the tower.
\begin{enumerate}
  \item The fermionic (ladder) modes at level $n{+}1$ split between the two
\end{enumerate}
   summands; those in $\mathcal{A}_{n+1}^{\mathrm{trans}}$ are the modes
   created by the jet prolongation that have no pre-image at level $n$.

---

Since $\tau_n$ is an idempotent $\emph{$-endomorphism of a $C^}$-algebra, the
splitting $\mathrm{im}\,\tau_n \oplus \ker\,\tau_n$ is an orthogonal direct
sum.

For item (1), let $\alpha \in \Gamma(J^1 E_{n+1})$ be a bosonic mode lifted
from level $n$: $[F_A, \alpha] = 0$ (\hyperref[rem:algebraic-filter]{Remark}).
Bosonic modes take values in the Cartan subalgebra
$\mathrm{U}(1) \subset \mathfrak{su}(2)$
(\hyperref[rem:weight-space-decomposition]{Remark}).  Since $\alpha$ lies in the
image of the ascent $\ell_n$, we have $\alpha = \ell_n(\beta)$ for some
$\beta \in \Gamma(J^1 E_n)$.  Then
$\tau_n(\alpha) = \ell_n(\mathfrak{d}_{n+1}(\ell_n(\beta)))
= \ell_n(\beta) = \alpha$,
using $\mathfrak{d}_{n+1} \circ \ell_n = \mathrm{id}$.  Hence
$\alpha \in \mathcal{A}_{n+1}^{\mathrm{desc}}$.

For item (2), the ladder modes $\omega^\pm \otimes E_\pm$ at level $n{+}1$
couple to the curvature $F_A$ through the covariant derivative $d_A$.  Those
that arise as lifts of level-$n$ modes lie in
$\mathcal{A}_{n+1}^{\mathrm{desc}}$ by the same argument as item (1).  The
remaining modes---genuinely new harmonic directions in $J^1 E_{n+1}$ with no
counterpart at level $n$---lie in $\mathcal{A}_{n+1}^{\mathrm{trans}}$ and
are annihilated by $\tau_n$: the descent $\mathfrak{d}_{n+1}$ maps them to
zero in the level-$n$ algebra, so $\tau_n = \ell_n \circ \mathfrak{d}_{n+1}$
sends them to zero.
\end{proposition}
\begin{proof}
Since $\tau_n$ is an idempotent $\emph{$-endomorphism of a $C^}$-algebra, the
splitting $\mathrm{im}\,\tau_n \oplus \ker\,\tau_n$ is an orthogonal direct
sum.

For item (1), let $\alpha \in \Gamma(J^1 E_{n+1})$ be a bosonic mode lifted
from level $n$: $[F_A, \alpha] = 0$ (\hyperref[rem:algebraic-filter]{Remark}).
Bosonic modes take values in the Cartan subalgebra
$\mathrm{U}(1) \subset \mathfrak{su}(2)$
(\hyperref[rem:weight-space-decomposition]{Remark}).  Since $\alpha$ lies in the
image of the ascent $\ell_n$, we have $\alpha = \ell_n(\beta)$ for some
$\beta \in \Gamma(J^1 E_n)$.  Then
$\tau_n(\alpha) = \ell_n(\mathfrak{d}_{n+1}(\ell_n(\beta)))
= \ell_n(\beta) = \alpha$,
using $\mathfrak{d}_{n+1} \circ \ell_n = \mathrm{id}$.  Hence
$\alpha \in \mathcal{A}_{n+1}^{\mathrm{desc}}$.

For item (2), the ladder modes $\omega^\pm \otimes E_\pm$ at level $n{+}1$
couple to the curvature $F_A$ through the covariant derivative $d_A$.  Those
that arise as lifts of level-$n$ modes lie in
$\mathcal{A}_{n+1}^{\mathrm{desc}}$ by the same argument as item (1).  The
remaining modes---genuinely new harmonic directions in $J^1 E_{n+1}$ with no
counterpart at level $n$---lie in $\mathcal{A}_{n+1}^{\mathrm{trans}}$ and
are annihilated by $\tau_n$: the descent $\mathfrak{d}_{n+1}$ maps them to
zero in the level-$n$ algebra, so $\tau_n = \ell_n \circ \mathfrak{d}_{n+1}$
sends them to zero.
\end{proof}

% ERROR: block "brane-boundary" not found
\begin{proposition}[Structure of the defining boundary]
  \label{prop:brane-boundary}
  \uses{def:brane-boundary, def:descent-involution, def:quantum-characteristic-classes, lem:descent}
The defining boundary $\partial A_n$
(\hyperref[def:brane-boundary]{Definition}) has the following properties:

\begin{enumerate}
  \item \emph{Codimension.}  In the archimedean realisation,
\end{enumerate}
   $\partial A_n$ corresponds to a closed subset of real codimension at
   least 2 in $E_n$.
\begin{enumerate}
  \item \emph{Exceptional object.}  At each degeneration locus of the
\end{enumerate}
   restriction $r_n$, the kernel $\ker(r_n|_S)$ is isomorphic to a
   simple object of $\mathcal{C}$; in the archimedean realisation
   this is $\mathbb{P}^1(\mathbb{C}) \cong S^2$, the exceptional divisor of the
   topological blow-up.
\begin{enumerate}
  \item \emph{Involution symmetry.}  The descent involution $\tau_n$
\end{enumerate}
   (\hyperref[def:descent-involution]{Definition}) restricts to an
   endomorphism of $\partial A_n$, whose fixed subobjects coincide
   with the monodromy fixed points of the quantum characteristic
   classes (\hyperref[def:quantum-characteristic-classes]{Definition}).
\begin{enumerate}
  \item \emph{Chirality interface (archimedean realisation).}
\end{enumerate}
   When $n = 1$, $\partial E_1$ separates the macroscopic spacetime
   coordinates from the internal chiral modes: it is the geometric
   boundary between the base manifold $M$ and the gauge fibre, across
   which the Reeb involution $\xi \mapsto -\xi$
   (\hyperref[rem:descartes-chirality]{Remark}) reverses orientation.

---

\emph{(1).} The descent map $\mathfrak{d}_{n+1}$ is a surjective
$*$-homomorphism of fibre algebras (\hyperref[lem:descent]{Lemma}).  Its kernel at a
generic point is zero because the ascent $\ell_n$ embeds $A$ as a direct
summand of $J^1(A)$ and the restriction $r_n$ is the left-inverse of this
embedding.  The kernel becomes non-trivial only where the direct-summand
splitting degenerates, which is a closed algebraic condition on the
Frobenius pairing $\varepsilon \circ m$.  Since $A$ is a Frobenius--Hopf
object in a rigid monoidal category, this degeneracy locus has real
codimension at least 2 by the standard codimension bound for degeneracy loci
of morphisms between vector bundles of equal rank.

\emph{(2).} At a blow-up centre $p \in \partial E_n$ the kernel of
$r_n|_p$ is the fibre of $\Omega^1(A)$ over $p$ modulo the image of
$\eta\colon \mathbf{1} \to A$.  Since
$\Omega^1(A) \hookrightarrow A \otimes A$ and $A$ has fibre
$\mathrm{SU}(2)$, the quotient is isomorphic to the projectivisation of
the $2$-dimensional representation, i.e.
$\mathbb{P}^1(\mathbb{C}) \cong S^2$.  This is the exceptional divisor introduced
by the topological blow-up that resolves the singularity at $p$.

\emph{(3).} The descent involution $\tau_n = \ell_n \circ \mathfrak{d}_{n+1}$
(\hyperref[def:descent-involution]{Definition}), acting on sections of
$J^1 E_{n+1}$, restricts to the boundary locus $\partial E_n$ as follows.
Since $\partial E_n$ is defined by $\ker(\mathfrak{d}_{n+1}|_p) \neq 0$,
and $\tau_n$ projects onto the image of $\ell_n$, the restriction
$\tau_n|_{\partial E_n}$ acts as an involution whose fixed points are
the sections at boundary points that are entirely in the descended
subalgebra---precisely the monodromy fixed-point condition that
defines the quantum characteristic classes
(\hyperref[def:quantum-characteristic-classes]{Definition}).

\emph{(4).} At $n = 1$ the base manifold is $M$ and the total space
$E_1$ fibres over $M$ with fibre $\mathrm{SU}(2)/\mathrm{U}(1) \cong S^2$.
The defining boundary $\partial E_1$ is the locus where the jet
restriction degenerates; by item (2) each fibre of $\partial E_1$ is
$\mathbb{P}^1(\mathbb{C})$, which separates the base (spacetime) coordinates from the
fibre (chiral) coordinates.  The Reeb vector field $\xi$ is vertical, so
its reversal $\xi \mapsto -\xi$ exchanges the two sides of $\partial E_1$:
one side carries left-handed fermions (positive Reeb eigenspace), the other
right-handed (negative eigenspace), in agreement with
\hyperref[rem:descartes-chirality]{Remark}.
\end{proposition}
\begin{proof}
\emph{(1).} The descent map $\mathfrak{d}_{n+1}$ is a surjective
$*$-homomorphism of fibre algebras (\hyperref[lem:descent]{Lemma}).  Its kernel at a
generic point is zero because the ascent $\ell_n$ embeds $A$ as a direct
summand of $J^1(A)$ and the restriction $r_n$ is the left-inverse of this
embedding.  The kernel becomes non-trivial only where the direct-summand
splitting degenerates, which is a closed algebraic condition on the
Frobenius pairing $\varepsilon \circ m$.  Since $A$ is a Frobenius--Hopf
object in a rigid monoidal category, this degeneracy locus has real
codimension at least 2 by the standard codimension bound for degeneracy loci
of morphisms between vector bundles of equal rank.

\emph{(2).} At a blow-up centre $p \in \partial E_n$ the kernel of
$r_n|_p$ is the fibre of $\Omega^1(A)$ over $p$ modulo the image of
$\eta\colon \mathbf{1} \to A$.  Since
$\Omega^1(A) \hookrightarrow A \otimes A$ and $A$ has fibre
$\mathrm{SU}(2)$, the quotient is isomorphic to the projectivisation of
the $2$-dimensional representation, i.e.
$\mathbb{P}^1(\mathbb{C}) \cong S^2$.  This is the exceptional divisor introduced
by the topological blow-up that resolves the singularity at $p$.

\emph{(3).} The descent involution $\tau_n = \ell_n \circ \mathfrak{d}_{n+1}$
(\hyperref[def:descent-involution]{Definition}), acting on sections of
$J^1 E_{n+1}$, restricts to the boundary locus $\partial E_n$ as follows.
Since $\partial E_n$ is defined by $\ker(\mathfrak{d}_{n+1}|_p) \neq 0$,
and $\tau_n$ projects onto the image of $\ell_n$, the restriction
$\tau_n|_{\partial E_n}$ acts as an involution whose fixed points are
the sections at boundary points that are entirely in the descended
subalgebra---precisely the monodromy fixed-point condition that
defines the quantum characteristic classes
(\hyperref[def:quantum-characteristic-classes]{Definition}).

\emph{(4).} At $n = 1$ the base manifold is $M$ and the total space
$E_1$ fibres over $M$ with fibre $\mathrm{SU}(2)/\mathrm{U}(1) \cong S^2$.
The defining boundary $\partial E_1$ is the locus where the jet
restriction degenerates; by item (2) each fibre of $\partial E_1$ is
$\mathbb{P}^1(\mathbb{C})$, which separates the base (spacetime) coordinates from the
fibre (chiral) coordinates.  The Reeb vector field $\xi$ is vertical, so
its reversal $\xi \mapsto -\xi$ exchanges the two sides of $\partial E_1$:
one side carries left-handed fermions (positive Reeb eigenspace), the other
right-handed (negative eigenspace), in agreement with
\hyperref[rem:descartes-chirality]{Remark}.
\end{proof}

\begin{remark}[Reduction at dimension 1]
  \label{rem:descent-reduction}
  \uses{def:q-reeb-vector-field, rem:algebraic-filter, rem:descent-reduction-zero}
When $\dim M_k = 1$ the geometric data of the archimedean quantum observable
universe collapse as follows.

\begin{enumerate}
  \item \emph{Connection.}  The $\mathrm{SU}(2)$ connection $A_k$ restricts to
\end{enumerate}
   a single dimension; its holonomy reduces to a $\mathrm{U}(1)$ phase
   around $S^1$.
\begin{enumerate}
  \item \emph{Bundle.}  The principal $\mathrm{SU}(2)$ bundle $E$ descends to a
\end{enumerate}
   bundle over the $1$-dimensional base whose fibre dynamics are locked
   into a single periodic flow.
\begin{enumerate}
  \item \emph{Reeb direction.}  The inverse of the Reeb vector field
\end{enumerate}
   (\hyperref[def:q-reeb-vector-field]{Definition}) becomes the only direction
   of motion; the horizontal distribution is trivial and no bosonic modes
   (\hyperref[rem:algebraic-filter]{Remark}) survive.

The descent need not terminate here; a further step carries $M_k$ to a
zero-dimensional base (\hyperref[rem:descent-reduction-zero]{Remark}).
\end{remark}

\begin{remark}[Physical implications of a $1$-dimensional base]
  \label{rem:descent-physical}
  \uses{def:mass-identity, rem:betti-numbers-brane}
At dimension 1 the physical content simplifies drastically.

\begin{enumerate}
  \item \emph{Primordial vortex.}  The vortex filaments of the mass identity
\end{enumerate}
   (\hyperref[def:mass-identity]{Definition}) are no longer embedded in a larger
   space; the entire manifold \emph{is} the filament.
\begin{enumerate}
  \item \emph{Mass as pure drag.}  When $\dim M_k = 1$ the surface
\end{enumerate}
   $\Sigma$ in the mass identity
   (\hyperref[def:mass-identity]{Definition}) degenerates to the
   $1$-manifold $M_k$ itself, and the $2$-dimensional nonabelian
   multiplicative integral reduces by the nonabelian Stokes
   theorem [yekutieli2015] to the holonomy of the connection
   $\nabla_A$ around $M_k$:
\[
m \;=\;
     \bigl\lVert\,
     \mathrm{Hol}_{\nabla_A}\!\bigl(M_k;\, \mathfrak{G}\bigr)
     \,\bigr\rVert,
\]
   so that mass represents the fundamental holonomy resistance of
   the $1$-dimensional flow against the $\mathrm{SU}(2)$ rotation.
\begin{enumerate}
  \item \emph{Lepton generations.}  Since the number of lepton generations is
\end{enumerate}
   bounded by the Euler characteristic and vertical degree
   (\hyperref[rem:betti-numbers-brane]{Remark}), a $1$-dimensional base supports
   only the most fundamental fermionic state; there is no room for
   higher-degree vertical harmonic forms.
\end{remark}

\begin{remark}[Reduction at dimension 0]
  \label{rem:descent-reduction-zero}
  \uses{def:q-reeb-vector-field}
When $\dim M_k = 0$ the base manifold is a discrete set of points and
the geometric data of the archimedean quantum observable universe
reduce to purely algebraic invariants at each point $p\in M_k$.

\begin{enumerate}
  \item \emph{Connection.}  The $\mathrm{SU}(2)$ connection $A_k$ has no
\end{enumerate}
   directions along which to parallel-transport; its curvature
   $F_A$ vanishes identically and the holonomy is trivial.
\begin{enumerate}
  \item \emph{Bundle.}  The principal $\mathrm{SU}(2)$ bundle $E$
\end{enumerate}
   restricts to its fibre $E_p$ over the point $p$; the bundle
   is determined entirely by this single fibre together with
   the group action.
\begin{enumerate}
  \item \emph{Substrate parameter.}  The substrate parameter
\end{enumerate}
   $q(p) = 1 + \lVert L\wedge F_A^{+}\rVert(p)$
   (\hyperref[rem:c-substrate-via-levi]{Remark}) assigns to each point a
   definite quantum constant $q\in[1,\infty)$, encoding the full
   content of the quantum observable universe at that point.
\begin{enumerate}
  \item \emph{Reeb direction.}  There is no tangent space, so the Reeb
\end{enumerate}
   vector field (\hyperref[def:q-reeb-vector-field]{Definition}) and the
   horizontal distribution are absent; only the internal
   $\mathrm{SU}(2)$ fibre data persist.

In this terminal stage every observable is determined by the fibre
$E_p$ and the value $q(p)$.  When $M_k$ consists of several points the
descent data decompose into independent summands, one per point; when
$M_k$ is connected---i.e. a single point---the entire quantum
observable universe is captured by the pair $(E_p, q(p))$.
\end{remark}

\begin{remark}[Physical implications of a $0$-dimensional base]
  \label{rem:descent-physical-zero}
  \uses{def:mass-identity, rem:betti-numbers-brane}
At dimension 0 all spatial extent has been removed and the physical
content is purely algebraic.

\begin{enumerate}
  \item \emph{Primordial state.}  The vortex filaments, flows and spatial
\end{enumerate}
   modes of higher-dimensional strata are absent; what remains is
   the initial datum of the fermionic content encoded in the
   fibre $E_p$.
\begin{enumerate}
  \item \emph{Mass.}  With no base manifold over which to integrate, the
\end{enumerate}
   nonabelian multiplicative integral of the mass identity
   (\hyperref[def:mass-identity]{Definition}) degenerates: the surface
   $\Sigma$ collapses to a point and the holonomy around $\gamma$
   reduces to a single fibre evaluation
   $m = \lVert L \wedge F^+\rVert(p)$.
   Mass at dimension 0 is an instantaneous, pointwise quantity
   rather than a distributed one.
\begin{enumerate}
  \item \emph{Lepton generations.}  All vertical harmonic forms
\end{enumerate}
   (\hyperref[rem:betti-numbers-brane]{Remark}) are concentrated at the
   point; the Euler characteristic $\chi(E_p)$ retains the count of
   fermionic degrees of freedom without any spatial multiplicity.
\end{remark}

\begin{remark}[Fixed-point mechanism: ascent and descent]
  \label{rem:descent-fixed-point}
  \uses{prop:jet-prolongation, lem:descent, def:descent-involution, def:brane-boundary, prop:descent-eigenspaces, rem:algebraic-filter, rem:descent-reduction, rem:descent-reduction-zero}
Each step down the brane tower removes a space-like dimension through the
$\mathrm{SU}(2)$ rotation and its associated fixed-point set.

\begin{enumerate}
  \item \emph{Ascent (generation).}  Moving from $M_k$ ($\dim = 0$) upward
\end{enumerate}
   lifts the connection $A$ to higher categorical structures via the jet
   prolongation (\hyperref[prop:jet-prolongation]{Proposition}).
\begin{enumerate}
  \item \emph{Descent (projection).}  Moving back down via the descent
\end{enumerate}
   homomorphism $\mathfrak{d}_n$ (\hyperref[lem:descent]{Lemma}) acts as an
   algebraic filter (\hyperref[rem:algebraic-filter]{Remark}) on the
   level-$(n{+}1)$ algebra.  The transient complement
   $\mathcal{A}_{n+1}^{\mathrm{trans}}$
   (\hyperref[prop:descent-eigenspaces]{Proposition})---the genuinely new
   modes at level $n{+}1$ with no pre-image at level $n$---is
   annihilated by the descent involution $\tau_n$.
   At dimension 1 only the longitudinal
   $\mathrm{U}(1)$ backbone remains
   (\hyperref[rem:descent-reduction]{Remark}), and at dimension 0 even
   that is collapsed to a single fibre datum
   (\hyperref[rem:descent-reduction-zero]{Remark}).
\begin{enumerate}
  \item \emph{Round-trip involution.}  The composition of ascent and descent
\end{enumerate}
   is the descent involution $\tau_n$
   (\hyperref[def:descent-involution]{Definition}), acting on
   $\Gamma(J^1 E_{n+1})$, whose image
   $\mathcal{A}_{n+1}^{\mathrm{desc}}$ is the stable core of the brane
   algebra at level $n{+}1$.  The defining boundary $\partial E_n$
   (\hyperref[def:brane-boundary]{Definition}) marks the geometric locus
   where the fibrewise descent map has a non-trivial kernel.
\end{remark}

\section{Topological Decomposition of Mass}
\label{sec:topological-mass}

\begin{theorem}[Topological Decomposition of Mass]
  \label{thm:topological-mass}
  \uses{def:quantum-observable-universe, def:recursive-brane-tower, def:mass-identity, def:descent-involution, def:brane-boundary, def:briggs-klein-chart, prop:descent-eigenspaces}
Let $(\mathcal{C}, A, \omega)$ be a quantum observable universe
(\hyperref[def:quantum-observable-universe]{Definition}); let $q$ denote the substrate parameter.  Here $\mathcal{C}$
is the rigid monoidal category, $A$ the Frobenius--Hopf object, and
$\omega$ the fibre functor.  Let
$\gamma \colon \mathbf{1} \to A$ be a morphism (a "closed cycle" in $A$,
\hyperref[def:mass-identity]{Definition}) whose categorical support is
generated by $N$ consecutive extensions along the brane tower
$\{A_n\}_{n \geq 0}$
(\hyperref[def:recursive-brane-tower]{Definition}), applied to a base-level
morphism $\gamma_0 \colon \mathbf{1} \to A_0$.

The mass $m(\gamma)$ (\hyperref[def:mass-identity]{Definition}) admits a
decomposition into descended and transient contributions via the
descent involution $\tau_n$
(\hyperref[def:descent-involution]{Definition}):

\begin{equation}\label{eq:topological-mass-decomp}
  m(\gamma)
  \;=\;
  \bigl\lVert\,
  \mathrm{Hol}_{\nabla_A}\!\bigl(\gamma_0;\, \mathfrak{G}\bigr)
  \,\bigr\rVert
  \;+\;
  \sum_{i=1}^{N}
  \bigl\lVert\,
  \mathrm{MI}\!\!\int_{E_i}
  \bigl(\nabla_A q,\; L \wedge F^+\bigr)
  \,\bigr\rVert,
\end{equation}

where:

\begin{enumerate}
  \item The first term is the holonomy of $\nabla_A$ around the unknot
\end{enumerate}
   core $\gamma_0$ in the twisted gerbe $\mathfrak{G}$, computed via the
   nonabelian Stokes theorem [yekutieli2015].  Categorically,
   this is the trace of the comultiplication $\Delta$ on $A$
   restricted to the descended subalgebra
   $\mathcal{A}^{\mathrm{desc}}$
   (\hyperref[prop:descent-eigenspaces]{Proposition}): it captures the
   baseline geometric drag of the Frobenius--Hopf object along the
   trivial core.
\begin{enumerate}
  \item Each summand is the nonabelian multiplicative integral of the
\end{enumerate}
   connection--curvature pair over the $i$-th exceptional divisor
   $E_i \cong \mathbb{P}^1(\mathbb{C})$ arising from the blow-up.
   Categorically, $E_i$ is the fibre of
   $\ker(\mathfrak{d}_{n+1}|_{p_i})$ at the $i$-th blow-up centre
   (\hyperref[def:brane-boundary]{Definition}), and the integral extracts
   the contribution of the transient complement
   $\mathcal{A}^{\mathrm{trans}}$ localised at $p_i$.  Evaluation
   of this integral via the spectral coordinates of the Briggs--Klein
   chart (\hyperref[def:briggs-klein-chart]{Definition}) yields the
   quantum dilogarithm $\Phi_q(z_i)$.

---

The proof proceeds in four steps: decomposition of the mass integral
using the descent involution, identification of the descended term with
holonomy, localisation of the transient terms to exceptional divisors,
and evaluation via the quantum dilogarithm.

\emph{Step 1: Eigenspace decomposition of the mass integrand.}
By the mass identity (\hyperref[def:mass-identity]{Definition}), the mass of
the vortex $\gamma$ is
\[
m(\gamma)
  \;=\;
  \bigl\lVert\,
  \mathrm{MI}\!\!\int_\Sigma
  \bigl(\nabla_A q,\; L \wedge F^+\bigr)
  \,\bigr\rVert,
\]
where $\Sigma$ is a surface with $\partial\Sigma = \gamma$.  The
integrand $(\nabla_A q,\; L \wedge F^+)$ is a section of the
connection--curvature bundle over $\Sigma$, and as such it is an element
of the section algebra $\Gamma(J^1 E_{n+1},\, \mathrm{SU}_{q_{n+1}}(2))$
at whatever level $n{+}1$ accommodates the vortex.

The descent involution $\tau_n$ (\hyperref[def:descent-involution]{Definition})
decomposes this section algebra into the direct sum
$\mathcal{A}_{n+1}^{\mathrm{desc}} \oplus \mathcal{A}_{n+1}^{\mathrm{trans}}$
(\hyperref[prop:descent-eigenspaces]{Proposition}).  Since $\tau_n$ is an
idempotent $*$-endomorphism, the Killing-form norm
$\lVert\cdot\rVert$ on $\mathfrak{su}(2)$ respects the orthogonal
decomposition:
\[
\lVert \alpha \rVert^2
  \;=\;
  \lVert \tau_n(\alpha) \rVert^2
  \;+\;
  \lVert \alpha - \tau_n(\alpha) \rVert^2
\]
for every section $\alpha$.  Applying this to the mass integrand and
using the additivity of the nonabelian multiplicative integral over
orthogonal direct sums [Proposition 5.3, yekutieli2015]:

\begin{equation}\label{eq:mi-decomp}
  \mathrm{MI}\!\!\int_\Sigma \alpha
  \;=\;
  \mathrm{MI}\!\!\int_\Sigma \tau_n(\alpha)
  \;\cdot\;
  \mathrm{MI}\!\!\int_\Sigma
  \bigl(\alpha - \tau_n(\alpha)\bigr),
\end{equation}

where the product is in the group $\mathrm{SU}_{q}(2)$.  Taking norms
and using the submultiplicativity of $\lVert\cdot\rVert$ (which is an
equality for the orthogonal decomposition of an idempotent):
\[
m(\gamma)
  \;=\;
  \bigl\lVert\,
  \mathrm{MI}\!\!\int_\Sigma \tau_n(\alpha)
  \,\bigr\rVert
  \;+\;
  \bigl\lVert\,
  \mathrm{MI}\!\!\int_\Sigma
  \bigl(\alpha - \tau_n(\alpha)\bigr)
  \,\bigr\rVert.
\]

\emph{Step 2: Descended term is the unknot-core holonomy.}
The projection $\tau_n(\alpha)$ lies in
$\mathcal{A}_{n+1}^{\mathrm{desc}} = \mathrm{im}\,\tau_n$, which is
isomorphic to the level-$n$ section algebra via the ascent embedding
$\ell_n$ (\hyperref[def:descent-involution]{Definition}).  The surface
$\Sigma$ with $\partial\Sigma = \gamma$ decomposes under the topological
blow-up as $\Sigma = \Sigma_0 \cup \bigcup_{i=1}^N D_i$, where
$\Sigma_0$ is a disc spanning the unknot core $\gamma_0$ and each
$D_i$ is the disc (exceptional divisor) inserted at the $i$-th blow-up.
On $\Sigma_0$ the transient component vanishes
($\alpha - \tau_n(\alpha) = 0$ on the unknot complement, since $\gamma_0$
has trivial topology and all sections there factor through level $n$).
Hence:
\[
\mathrm{MI}\!\!\int_\Sigma \tau_n(\alpha)
  \;=\;
  \mathrm{MI}\!\!\int_{\Sigma_0} \alpha
  \;=\;
  \mathrm{Hol}_{\nabla_A}\!\bigl(\gamma_0;\, \mathfrak{G}\bigr),
\]
where the second equality is the nonabelian Stokes
theorem [yekutieli2015] applied to the disc $\Sigma_0$.

\emph{Step 3: Transient terms localise to exceptional divisors.}
The transient component $\alpha - \tau_n(\alpha)$ is annihilated by
$\tau_n$: it lies in
$\mathcal{A}_{n+1}^{\mathrm{trans}} = \ker\,\tau_n$
(\hyperref[prop:descent-eigenspaces]{Proposition}).  By the structure of the
defining boundary (\hyperref[prop:brane-boundary]{Proposition}), the support of
a transient section is contained in a neighbourhood of the boundary
$\partial E_n$, which has codimension $\geq 2$ in $E_n$ (item (1)) and
whose fibre at each blow-up centre $p_i$ is the exceptional divisor
$E_i \cong \mathbb{P}^1(\mathbb{C})$ (item (2)).

On the surface $\Sigma$, the transient integrand is therefore supported
on the discs $D_i$ centred at the blow-up points.  Since the $D_i$ are
disjoint (each blow-up introduces a new connected component of
$\partial E_n$), the multiplicative integral splits as a product:
\[
\mathrm{MI}\!\!\int_\Sigma
  \bigl(\alpha - \tau_n(\alpha)\bigr)
  \;=\;
  \prod_{i=1}^{N}\;
  \mathrm{MI}\!\!\int_{D_i}
  \bigl(\nabla_A q,\; L \wedge F^+\bigr).
\]
Since the $D_i$ are geometrically disjoint and the group elements commute
up to the central band $\mathcal{O}_M^{\times}$ of the gerbe
(the transient modes are localised at distinct points, and the central
extension ensures that the group-valued integrals commute modulo the
band), the norm of the product equals the sum of the norms:
\[
\bigl\lVert\,
  \prod_{i=1}^{N}\;
  \mathrm{MI}\!\!\int_{D_i}
  \bigl(\nabla_A q,\; L \wedge F^+\bigr)
  \,\bigr\rVert
  \;=\;
  \sum_{i=1}^{N}
  \bigl\lVert\,
  \mathrm{MI}\!\!\int_{E_i}
  \bigl(\nabla_A q,\; L \wedge F^+\bigr)
  \,\bigr\rVert,
\]
where we replace $D_i$ by $E_i$ since the integrand is supported on the
exceptional fibre $E_i \subset D_i$.

\emph{Step 4: Quantum dilogarithm evaluation.}
Combining Steps 1--3:
\[
m(\gamma)
  \;=\;
  \bigl\lVert\,
  \mathrm{Hol}_{\nabla_A}\!\bigl(\gamma_0;\, \mathfrak{G}\bigr)
  \,\bigr\rVert
  \;+\;
  \sum_{i=1}^{N}
  \bigl\lVert\,
  \mathrm{MI}\!\!\int_{E_i}
  \bigl(\nabla_A q,\; L \wedge F^+\bigr)
  \,\bigr\rVert.
\]
This is the claimed decomposition \hyperref[eq:topological-mass-decomp]{eq:topological-mass-decomp}.
The identification of each summand with the quantum dilogarithm
$\Phi_q(z_i)$ follows from the Corollary below: integration of the
$q$-deformed expansion flow over the resolved defect $E_i$ in the
spectral coordinates of the Briggs--Klein chart yields
$\Phi_q(z_i) = \int_0^{z_i} \ln_q(1+t)/t\, dt$.
\end{theorem}
\begin{proof}
The proof proceeds in four steps: decomposition of the mass integral
using the descent involution, identification of the descended term with
holonomy, localisation of the transient terms to exceptional divisors,
and evaluation via the quantum dilogarithm.

\emph{Step 1: Eigenspace decomposition of the mass integrand.}
By the mass identity (\hyperref[def:mass-identity]{Definition}), the mass of
the vortex $\gamma$ is
\[
m(\gamma)
  \;=\;
  \bigl\lVert\,
  \mathrm{MI}\!\!\int_\Sigma
  \bigl(\nabla_A q,\; L \wedge F^+\bigr)
  \,\bigr\rVert,
\]
where $\Sigma$ is a surface with $\partial\Sigma = \gamma$.  The
integrand $(\nabla_A q,\; L \wedge F^+)$ is a section of the
connection--curvature bundle over $\Sigma$, and as such it is an element
of the section algebra $\Gamma(J^1 E_{n+1},\, \mathrm{SU}_{q_{n+1}}(2))$
at whatever level $n{+}1$ accommodates the vortex.

The descent involution $\tau_n$ (\hyperref[def:descent-involution]{Definition})
decomposes this section algebra into the direct sum
$\mathcal{A}_{n+1}^{\mathrm{desc}} \oplus \mathcal{A}_{n+1}^{\mathrm{trans}}$
(\hyperref[prop:descent-eigenspaces]{Proposition}).  Since $\tau_n$ is an
idempotent $*$-endomorphism, the Killing-form norm
$\lVert\cdot\rVert$ on $\mathfrak{su}(2)$ respects the orthogonal
decomposition:
\[
\lVert \alpha \rVert^2
  \;=\;
  \lVert \tau_n(\alpha) \rVert^2
  \;+\;
  \lVert \alpha - \tau_n(\alpha) \rVert^2
\]
for every section $\alpha$.  Applying this to the mass integrand and
using the additivity of the nonabelian multiplicative integral over
orthogonal direct sums [Proposition 5.3, yekutieli2015]:

\begin{equation}\label{eq:mi-decomp}
  \mathrm{MI}\!\!\int_\Sigma \alpha
  \;=\;
  \mathrm{MI}\!\!\int_\Sigma \tau_n(\alpha)
  \;\cdot\;
  \mathrm{MI}\!\!\int_\Sigma
  \bigl(\alpha - \tau_n(\alpha)\bigr),
\end{equation}

where the product is in the group $\mathrm{SU}_{q}(2)$.  Taking norms
and using the submultiplicativity of $\lVert\cdot\rVert$ (which is an
equality for the orthogonal decomposition of an idempotent):
\[
m(\gamma)
  \;=\;
  \bigl\lVert\,
  \mathrm{MI}\!\!\int_\Sigma \tau_n(\alpha)
  \,\bigr\rVert
  \;+\;
  \bigl\lVert\,
  \mathrm{MI}\!\!\int_\Sigma
  \bigl(\alpha - \tau_n(\alpha)\bigr)
  \,\bigr\rVert.
\]

\emph{Step 2: Descended term is the unknot-core holonomy.}
The projection $\tau_n(\alpha)$ lies in
$\mathcal{A}_{n+1}^{\mathrm{desc}} = \mathrm{im}\,\tau_n$, which is
isomorphic to the level-$n$ section algebra via the ascent embedding
$\ell_n$ (\hyperref[def:descent-involution]{Definition}).  The surface
$\Sigma$ with $\partial\Sigma = \gamma$ decomposes under the topological
blow-up as $\Sigma = \Sigma_0 \cup \bigcup_{i=1}^N D_i$, where
$\Sigma_0$ is a disc spanning the unknot core $\gamma_0$ and each
$D_i$ is the disc (exceptional divisor) inserted at the $i$-th blow-up.
On $\Sigma_0$ the transient component vanishes
($\alpha - \tau_n(\alpha) = 0$ on the unknot complement, since $\gamma_0$
has trivial topology and all sections there factor through level $n$).
Hence:
\[
\mathrm{MI}\!\!\int_\Sigma \tau_n(\alpha)
  \;=\;
  \mathrm{MI}\!\!\int_{\Sigma_0} \alpha
  \;=\;
  \mathrm{Hol}_{\nabla_A}\!\bigl(\gamma_0;\, \mathfrak{G}\bigr),
\]
where the second equality is the nonabelian Stokes
theorem [yekutieli2015] applied to the disc $\Sigma_0$.

\emph{Step 3: Transient terms localise to exceptional divisors.}
The transient component $\alpha - \tau_n(\alpha)$ is annihilated by
$\tau_n$: it lies in
$\mathcal{A}_{n+1}^{\mathrm{trans}} = \ker\,\tau_n$
(\hyperref[prop:descent-eigenspaces]{Proposition}).  By the structure of the
defining boundary (\hyperref[prop:brane-boundary]{Proposition}), the support of
a transient section is contained in a neighbourhood of the boundary
$\partial E_n$, which has codimension $\geq 2$ in $E_n$ (item (1)) and
whose fibre at each blow-up centre $p_i$ is the exceptional divisor
$E_i \cong \mathbb{P}^1(\mathbb{C})$ (item (2)).

On the surface $\Sigma$, the transient integrand is therefore supported
on the discs $D_i$ centred at the blow-up points.  Since the $D_i$ are
disjoint (each blow-up introduces a new connected component of
$\partial E_n$), the multiplicative integral splits as a product:
\[
\mathrm{MI}\!\!\int_\Sigma
  \bigl(\alpha - \tau_n(\alpha)\bigr)
  \;=\;
  \prod_{i=1}^{N}\;
  \mathrm{MI}\!\!\int_{D_i}
  \bigl(\nabla_A q,\; L \wedge F^+\bigr).
\]
Since the $D_i$ are geometrically disjoint and the group elements commute
up to the central band $\mathcal{O}_M^{\times}$ of the gerbe
(the transient modes are localised at distinct points, and the central
extension ensures that the group-valued integrals commute modulo the
band), the norm of the product equals the sum of the norms:
\[
\bigl\lVert\,
  \prod_{i=1}^{N}\;
  \mathrm{MI}\!\!\int_{D_i}
  \bigl(\nabla_A q,\; L \wedge F^+\bigr)
  \,\bigr\rVert
  \;=\;
  \sum_{i=1}^{N}
  \bigl\lVert\,
  \mathrm{MI}\!\!\int_{E_i}
  \bigl(\nabla_A q,\; L \wedge F^+\bigr)
  \,\bigr\rVert,
\]
where we replace $D_i$ by $E_i$ since the integrand is supported on the
exceptional fibre $E_i \subset D_i$.

\emph{Step 4: Quantum dilogarithm evaluation.}
Combining Steps 1--3:
\[
m(\gamma)
  \;=\;
  \bigl\lVert\,
  \mathrm{Hol}_{\nabla_A}\!\bigl(\gamma_0;\, \mathfrak{G}\bigr)
  \,\bigr\rVert
  \;+\;
  \sum_{i=1}^{N}
  \bigl\lVert\,
  \mathrm{MI}\!\!\int_{E_i}
  \bigl(\nabla_A q,\; L \wedge F^+\bigr)
  \,\bigr\rVert.
\]
This is the claimed decomposition \hyperref[eq:topological-mass-decomp]{eq:topological-mass-decomp}.
The identification of each summand with the quantum dilogarithm
$\Phi_q(z_i)$ follows from the Corollary below: integration of the
$q$-deformed expansion flow over the resolved defect $E_i$ in the
spectral coordinates of the Briggs--Klein chart yields
$\Phi_q(z_i) = \int_0^{z_i} \ln_q(1+t)/t\, dt$.
\end{proof}

\begin{corollary}[Quantum Dilogarithm as Localized Torsion Volume]
  \label{cor:quantum-dilogarithm}
  \uses{thm:topological-mass}
The localized topological mass contribution $\Phi_q(z_i)$ at the $i$-th exceptional divisor $E_i$ is exactly the geometric integration of the 1-dimensional $q$-deformed expansion flow over the resolved defect. Integrating the spectral uncertainty flow $\ln_q$ along the radial coordinate of the Briggs-Klein affine chart yields the quantum dilogarithm:

\begin{equation}
\Phi_q(z_i) = \int_0^{z_i} \frac{\ln_q(1+t)}{t} dt = \sum_{n=1}^\infty (-1)^{n+1} \frac{z_i^n}{n [n]_q}
\end{equation}

where $\ln_q(1+t) = \sum_{n=1}^\infty (-1)^{n+1} \frac{t^n}{[n]_q}$ governs the continuous Reeb flow. The integration introduces the factor of $n$ in the denominator, elevating the 1-dimensional path propagation to a geometric volume measure trapping the topological charge.
\end{corollary}

\begin{remark}[Quantum dilogarithm integrand as $q$-deformed Böttcher potential density]
  \label{rem:qdlog-julia}
  \uses{def:substrate-parameter, def:q-reeb-flow, rem:reeb-exceptional-divisor}
The integrand $\ln_q(1+t)/t$ in the quantum dilogarithm formula above has the same formal
structure as the \emph{information-flow density} that appears in the computation of
filled Julia sets via the Bottcher escape potential.

For the quadratic family $f_c(z)=z^2+c$, the external Green's function (escape potential)
admits the convergent series

\begin{equation}\label{eq:bottcher-series}
  G_c(z)
  \;=\;
  \sum_{k=0}^{\infty}
    \frac{1}{2^{k+1}}\,
    \ln\!\left(1+\frac{c}{f_c^k(z)^2}\right).
\end{equation}

Writing $t_k = c/f_c^k(z)^2$, each summand is $\frac{1}{2^{k+1}}\ln(1+t_k)$: the
function $\ln(1+t_k)/t_k$ (evaluated at $t_k$) is exactly the classical ($q=1$) limit
of the quantum dilogarithm integrand.  The sum over $k$ weighted by $2^{-(k+1)}$
accumulates the \emph{information flow} along the orbit of $f_c$, with the weight
$2^{-k}$ reflecting the degree-$2$ expansion at each step.

In the QOU framework the role of the dynamical parameter $c$ is played by the
substrate parameter $q=q(z)$ (\hyperref[def:substrate-parameter]{Definition}), which varies
continuously with the spacetime coordinate $z$.  The $q$-deformed Reeb flow
(\hyperref[def:q-reeb-flow]{Definition}) on the exceptional divisor $E_i \cong \mathbb{P}^1(\mathbb{C})$
takes the place of the orbit $\{f_c^k(z)\}_{k\geq 0}$, and the path from $0$ to $z_i$
in the quantum dilogarithm integral accumulates information along that flow exactly as
\hyperref[eq:bottcher-series]{eq:bottcher-series} accumulates it along the Julia-set orbit.  The
$q$-deformation of the logarithm encodes the departure from classical ($q=1$)
dynamics: as $q\to 1$ the flow trivialises and $\Phi_q(z_i)\to\mathrm{Li}_2(z_i)$,
the classical dilogarithm; as $q\to\infty$ the quantum-number denominators $[n]_q$
suppress high-order terms, concentrating the topological charge near the blow-up centre.

In this identification the \emph{Julia set boundary} of the classical problem---the
locus where $G_c$ transitions from zero to positive---is the analogue of the boundary
of the exceptional divisor $\partial E_i$, where the $q$-deformed torsion charge changes
regime (\hyperref[rem:reeb-exceptional-divisor]{Remark}).  The substrate parameter $q(z)$
thereby acts as a \emph{spacetime-dependent dynamical parameter}, continuously
deforming the Julia set geometry as one moves through the QOU.
\end{remark}

\begin{remark}[Julia set on $S^1$ as photon scattering at the Big Bang]
  \label{rem:julia-s1-photon-scatter}
  \uses{def:q-reeb-flow, def:boson, rem:reeb-exceptional-divisor}
The exceptional divisor $E_i\cong\mathbb{P}^1(\mathbb{C})$ contains the equatorial circle
$S^1\subset E_i$.  The $q$-deformed Reeb flow (\hyperref[def:q-reeb-flow]{Definition})
restricts to a degree-$2$ endomorphism $f_q\colon S^1\to S^1$, whose
\emph{Julia set} $J(f_q)\subset S^1$ is the closure of repelling periodic
points---the locus of non-equicontinuous dynamics on the circle.

Bosons are purely horizontal $q$-harmonic forms (\hyperref[def:boson]{Definition});
photons are the empty-knot background flux on $E_i$.  The three regions of the
circle map identify as:

\begin{itemize}
  \item \emph{Filled Julia set} $\mathcal{K}(f_q)\cap S^1$: photons trapped in
\end{itemize}
  the $q$-torsion plasma---the opaque pre-recombination regime.
\begin{itemize}
  \item \emph{Julia set boundary} $\partial\mathcal{K}(f_q)\cap S^1$: the
\end{itemize}
  surface of last scattering, where photons transition from trapped to
  free-streaming.
\begin{itemize}
  \item \emph{Escape region} $S^1\setminus\mathcal{K}(f_q)$: freely propagating
\end{itemize}
  photons in the transparent universe.

At the Big Bang ($q\to\infty$, \hyperref[rem:reeb-exceptional-divisor]{Remark}), the
denominators $[n]_q$ suppress all escape and $J(f_q)=S^1$: total opacity.
As $q\to 1$ (heat death), $f_q$ trivialises, $J(f_q)$ collapses to a Cantor
set of measure zero, and the universe becomes transparent.
\end{remark}
